\chapter{Chapter07 — Algebraic PCPP}
%%% generated by the blueprint pipeline from FABL.Chapter07.AlgebraicPCPP
\section{Overview}

Declarations of the Lean module \texttt{FABL.Chapter07.AlgebraicPCPP}.

\section{Declarations}

\begin{definition}[Binary subspace splitting data]
\lean{FABL.BinarySubspaceSplit}
\statementsource{Analysis-of-Boolean-Functions-by-Ryan-ODonnell}{
  pdf-ffbe9c7b3306-p072-b004,
  pdf-ffbe9c7b3306-p073-b002,
  pdf-ffbe9c7b3306-p197-b009
}
\label{FABL.BinarySubspaceSplit}
\leanok
\uses{FABL.F₂Cube, FABL.𝔽₂}
Let $H$ be a subspace of the $n$-dimensional vector space over $bF_2$. A binary
subspace splitting consists of a finite enumeration of $H$, a natural number $d$, and a
linear isomorphism from $H	imesbF_2^d$ to $bF_2^n$ such that every $(h,0)$ maps to the
original vector $h$ in $bF_2^n$.
\end{definition}

\begin{definition}[Existence of binary subspace splitting data]
\lean{FABL.binarySubspaceSplit}
\statementsource{Analysis-of-Boolean-Functions-by-Ryan-ODonnell}{
  pdf-ffbe9c7b3306-p072-b004,
  pdf-ffbe9c7b3306-p073-b002,
  pdf-ffbe9c7b3306-p197-b009
}
\label{FABL.binarySubspaceSplit}
\leanok
\uses{FABL.BinarySubspaceSplit, FABL.F₂Cube, FABL.𝔽₂}
For every subspace $H$ of $bF_2^n$, there is a binary subspace splitting: a natural
number $d$ and a linear isomorphism $H	imesbF_2^d	obF_2^n$ that sends $(h,0)$ to $h$
for every $h\in H$.
\end{definition}

\begin{definition}[Point in split coordinates]
\lean{FABL.BinarySubspaceSplit.point}
\statementsource{Analysis-of-Boolean-Functions-by-Ryan-ODonnell}{pdf-ffbe9c7b3306-p197-b009}
\label{FABL.BinarySubspaceSplit.point}
\leanok
\uses{FABL.BinarySubspaceSplit, FABL.F₂Cube, FABL.𝔽₂}
Given a binary subspace splitting of $H\subseteqbF_2^n$, its point associated to $h\in
H$ and $x\inbF_2^d$ is the image of $(h,x)$ under the splitting isomorphism.
\end{definition}

\begin{definition}[Character from split coordinates]
\lean{FABL.BinarySubspaceSplit.character}
\statementsource{Analysis-of-Boolean-Functions-by-Ryan-ODonnell}{
  pdf-ffbe9c7b3306-p071-b009,
  pdf-ffbe9c7b3306-p197-b009
}
\label{FABL.BinarySubspaceSplit.character}
\leanok
\uses{FABL.BinarySubspaceSplit, FABL.F₂Cube, FABL.Sign, FABL.f₂DotProduct, FABL.signEncode, FABL.𝔽₂}
Given a binary subspace splitting of $H\subseteqbF_2^n$ and a vector
$gamma\inbF_2^d$, define a sign-valued function on $bF_2^n$ by sending $z$ to
$(-1)^{langlegamma,yrangle}$, where $y$ is the second coordinate of the inverse image
of $z$ under the splitting isomorphism and the inner product is over $bF_2$.
\end{definition}

\begin{definition}[Split character]
\lean{FABL.IsSplitCharacter}
\statementsource{Analysis-of-Boolean-Functions-by-Ryan-ODonnell}{
  pdf-ffbe9c7b3306-p071-b009,
  pdf-ffbe9c7b3306-p072-b004,
  pdf-ffbe9c7b3306-p197-b009
}
\label{FABL.IsSplitCharacter}
\leanok
\uses{FABL.BinarySubspaceSplit, FABL.BinarySubspaceSplit.character, FABL.F₂Cube, FABL.Sign, FABL.𝔽₂}
Given a binary subspace splitting of $H\subseteqbF_2^n$, a sign-valued function on
$bF_2^n$ is a split character if it equals the character associated to some vector
$gamma\inbF_2^d$ by the splitting.
\end{definition}

\begin{definition}[Frequency functional]
\lean{FABL.BinarySubspaceSplit.frequencyFunctional}
\statementsource{Analysis-of-Boolean-Functions-by-Ryan-ODonnell}{pdf-ffbe9c7b3306-p197-b009}
\label{FABL.BinarySubspaceSplit.frequencyFunctional}
\leanok
\uses{FABL.BinarySubspaceSplit, FABL.F₂Cube, FABL.𝔽₂}
Given a binary subspace splitting of $H\subseteqbF_2^n$ and $gamma\inbF_2^d$, the
associated frequency functional sends $z\inbF_2^n$ to the $bF_2$ dot product of
$gamma$ with the second coordinate of the inverse image of $z$ under the splitting
isomorphism.
\end{definition}

\begin{definition}[Vector representing a frequency]
\lean{FABL.BinarySubspaceSplit.liftFrequency}
\statementsource{Analysis-of-Boolean-Functions-by-Ryan-ODonnell}{pdf-ffbe9c7b3306-p197-b009}
\label{FABL.BinarySubspaceSplit.liftFrequency}
\leanok
\uses{FABL.BinarySubspaceSplit, FABL.BinarySubspaceSplit.frequencyFunctional, FABL.F₂Cube, FABL.𝔽₂}
Given a binary subspace splitting of $H\subseteqbF_2^n$ and $gamma\inbF_2^d$, the
lifted frequency is the unique vector in $bF_2^n$ whose dot product with each
$z\inbF_2^n$ equals the frequency functional associated to $gamma$ evaluated at $z$.
\end{definition}

\begin{theorem}[Character and lifted frequency]
\lean{FABL.BinarySubspaceSplit.character_eq_liftFrequency}
\label{FABL.BinarySubspaceSplit.character_eq_liftFrequency}
\leanok
\uses{FABL.BinarySubspaceSplit, FABL.BinarySubspaceSplit.character, FABL.BinarySubspaceSplit.frequencyFunctional, FABL.BinarySubspaceSplit.liftFrequency, FABL.F₂Cube, FABL.f₂DotProduct, FABL.signEncode, FABL.𝔽₂}
\difficulty{3}
Let $H$ be a subspace of the binary cube $\br_2^n$, and let a binary subspace split
identify that cube linearly with $H$ times a complementary binary cube. For each
complementary frequency $\gamma$, the character that evaluates the dot product of
$\gamma$ with the complementary coordinate of a point $z$ equals the sign encoding of
the dot product of the lifted frequency with $z$.
\end{theorem}

\begin{theorem}[Lifted frequency is perpendicular]
\lean{FABL.BinarySubspaceSplit.liftFrequency_mem_perpendicularSubspace}
\label{FABL.BinarySubspaceSplit.liftFrequency_mem_perpendicularSubspace}
\leanok
\uses{FABL.BinarySubspaceSplit, FABL.BinarySubspaceSplit.frequencyFunctional, FABL.BinarySubspaceSplit.liftFrequency, FABL.F₂Cube, FABL.f₂DotProduct, FABL.mem_perpendicularSubspace_iff, FABL.perpendicularSubspace, FABL.𝔽₂}
\difficulty{4}
Let $H$ be a subspace of the binary cube $\bbr_2^n$, and let a binary subspace split
identify that cube linearly with $H$ times a complementary binary cube. For every
frequency $\gamma$ in the complementary cube, its lift to the original cube belongs to
the perpendicular subspace of $H$.
\end{theorem}

\begin{definition}[Restriction of the dot-product functional]
\lean{FABL.BinarySubspaceSplit.restrictionFunctional}
\statementsource{Analysis-of-Boolean-Functions-by-Ryan-ODonnell}{pdf-ffbe9c7b3306-p197-b009}
\label{FABL.BinarySubspaceSplit.restrictionFunctional}
\leanok
\uses{FABL.BinarySubspaceSplit, FABL.F₂Cube, FABL.𝔽₂}
Let $n$ be a natural number, let $H$ be a subspace of the binary cube $[𝔽_2]^n$, and
let a split identify $H 	imes [𝔽_2]^d$ linearly with $[𝔽_2]^n$, carrying $(h,0)$ to
$h$ for every $h [𝔽_2] H$. For $[𝔽_2]$-valued vectors $[η] [𝔽_2]^n$, the restriction
functional associated with this split sends $y [𝔽_2]^d$ to the dot product of $[η]$
with the image of $(0,y)$ under the identification.
\end{definition}

\begin{definition}[Restricted frequency]
\lean{FABL.BinarySubspaceSplit.restrictFrequency}
\statementsource{Analysis-of-Boolean-Functions-by-Ryan-ODonnell}{pdf-ffbe9c7b3306-p197-b009}
\label{FABL.BinarySubspaceSplit.restrictFrequency}
\leanok
\uses{FABL.BinarySubspaceSplit, FABL.BinarySubspaceSplit.restrictionFunctional, FABL.F₂Cube, FABL.𝔽₂}
Let $H$ be a subspace of the binary cube $br^n$ over $br/2br$, let a split of $H$
identify $H$ together with a complementary binary cube of dimension $d$ with the full
cube, and let $\eta$ be a frequency in the full cube. The restricted frequency is
the vector in the complementary cube whose dot-product functional corresponds, under
this identification, to the restriction functional of $\eta$.
\end{definition}

\begin{theorem}[Dot product of the restricted frequency]
\lean{FABL.BinarySubspaceSplit.dotProduct_restrictFrequency}
\label{FABL.BinarySubspaceSplit.dotProduct_restrictFrequency}
\leanok
\uses{FABL.BinarySubspaceSplit, FABL.BinarySubspaceSplit.restrictFrequency, FABL.F₂Cube, FABL.f₂DotProduct, FABL.𝔽₂}
\difficulty{3}
Let $H$ be a subspace of the binary cube $\bbr_2^n$, with a binary subspace split
identifying that cube linearly with $H$ times a complementary binary cube. For every
vector $\eta$ in the original cube and every complementary vector $a$, the dot product
of the frequency obtained by restricting $\eta$ to the complementary factor with $a$
equals the dot product of $\eta$ with the image of $(0,a)$ under the splitting
isomorphism.
\end{theorem}

\begin{theorem}[Character from restricted frequency]
\lean{FABL.BinarySubspaceSplit.character_restrictFrequency}
\label{FABL.BinarySubspaceSplit.character_restrictFrequency}
\leanok
\uses{FABL.BinarySubspaceSplit, FABL.BinarySubspaceSplit.character, FABL.BinarySubspaceSplit.dotProduct_restrictFrequency, FABL.BinarySubspaceSplit.restrictFrequency, FABL.F₂Cube, FABL.f₂DotProduct, FABL.mem_perpendicularSubspace_iff, FABL.perpendicularSubspace, FABL.signEncode, FABL.𝔽₂}
\difficulty{4}
Let $H$ be a subspace of the binary cube $\bbr_2^n$, and let a binary subspace split
identify that cube linearly with $H$ times a complementary binary cube. If $\eta$
belongs to the perpendicular subspace of $H$, then the character associated to the
frequency obtained by restricting $\eta$ to the complementary factor is the sign
encoding of the dot product of $\eta$ with the point $z$, for every $z$ in the original
cube.
\end{theorem}

\begin{definition}[Character from the perpendicular subspace]
\lean{FABL.IsPerpendicularCharacter}
\statementsource{Analysis-of-Boolean-Functions-by-Ryan-ODonnell}{
  pdf-ffbe9c7b3306-p072-b004,
  pdf-ffbe9c7b3306-p197-b009
}
\label{FABL.IsPerpendicularCharacter}
\leanok
\uses{FABL.F₂Cube, FABL.Sign, FABL.f₂DotProduct, FABL.perpendicularSubspace, FABL.signEncode, FABL.𝔽₂}
Let $H$ be a subspace of the binary cube $br^n$ over the field with two elements, and
let $f$ map the binary cube to the signs $\{\pm1\}$. Then $f$ is a character arising
from the perpendicular subspace of $H$ precisely when there is some $\eta$ in that
perpendicular subspace such that, for every $z$, $f(z)$ is the sign encoding of the
binary dot product of $\eta$ and $z$.
\end{definition}

\begin{theorem}[Split and perpendicular characters]
\lean{FABL.isSplitCharacter_iff_isPerpendicularCharacter}
\label{FABL.isSplitCharacter_iff_isPerpendicularCharacter}
\leanok
\uses{FABL.BinarySubspaceSplit, FABL.BinarySubspaceSplit.character_eq_liftFrequency, FABL.BinarySubspaceSplit.character_restrictFrequency, FABL.BinarySubspaceSplit.liftFrequency, FABL.BinarySubspaceSplit.liftFrequency_mem_perpendicularSubspace, FABL.BinarySubspaceSplit.restrictFrequency, FABL.F₂Cube, FABL.IsPerpendicularCharacter, FABL.IsSplitCharacter, FABL.Sign, FABL.𝔽₂}
\difficulty{4}
Let $n$ be a nonnegative integer, let $H$ be a linear subspace of $\mathbb F_2^n$, and
choose a linear isomorphism from $H\times\mathbb F_2^d$ to $\mathbb F_2^n$ for some
dimension $d$, whose restriction to the first factor agrees with the inclusion of $H$. A
sign character obtained by taking a dot product with a frequency in the complementary
coordinates is exactly a character whose frequency lies in the perpendicular subspace
$H^\perp$; equivalently, a function is representable in the chosen split-character form
if and only if it is representable as $z\mapsto(-1)^{\eta\cdot z}$ for some $\eta\in
H^\perp$.
\end{theorem}

\begin{definition}[Acceptance of a split character transcript]
\lean{FABL.splitCharacterTranscriptAccept}
\statementsource{Analysis-of-Boolean-Functions-by-Ryan-ODonnell}{pdf-ffbe9c7b3306-p197-b009}
\label{FABL.splitCharacterTranscriptAccept}
\leanok
\uses{FABL.Sign}
Given three optional signs $v_0,v_1,v_2$, accept exactly when all three are present and
the product of the first two equals the third; otherwise reject.
\end{definition}

\begin{definition}[Split character seed space]
\lean{FABL.SplitCharacterSeed}
\statementsource{Analysis-of-Boolean-Functions-by-Ryan-ODonnell}{pdf-ffbe9c7b3306-p197-b009}
\label{FABL.SplitCharacterSeed}
\leanok
\uses{FABL.BinarySubspaceSplit, FABL.F₂Cube, FABL.𝔽₂}
Let $H$ be a subspace of the binary cube $br_2^n$, and let a binary subspace split of
$H$ have complementary dimension $d$. The split character seed space is
$(br_2^d\times\bbr_2^d)\times H$.
\end{definition}

\begin{definition}[Split character function tester]
\lean{FABL.splitCharacterFunctionTester}
\statementsource{Analysis-of-Boolean-Functions-by-Ryan-ODonnell}{pdf-ffbe9c7b3306-p197-b009}
\label{FABL.splitCharacterFunctionTester}
\leanok
\uses{FABL.BinarySubspaceSplit, FABL.BinarySubspaceSplit.point, FABL.F₂Cube, FABL.NonadaptiveFunctionTester, FABL.Sign, FABL.SplitCharacterSeed, FABL.splitCharacterTranscriptAccept, FABL.uniformPMF, FABL.𝔽₂}
Given a binary subspace $H$ of the cube ${0,1}^n$ together with a linear decomposition
of the cube into $H$ and a complementary factor, this defines a three-query nonadaptive
randomized tester with a uniformly sampled seed. On seed $(h,u,v)$, it queries the three
points obtained by applying the split's point map to $(u,h)$, $(0,v)$, and $(0,u+v)$,
respectively, and accepts according to the split-character transcript predicate.
\end{definition}

\begin{theorem}[Perfect Completeness of the Split Character Tester]
\lean{FABL.splitCharacterFunctionTester_complete}
\label{FABL.splitCharacterFunctionTester_complete}
\leanok
\uses{FABL.BinarySubspaceSplit, FABL.BinarySubspaceSplit.character, FABL.BinarySubspaceSplit.point, FABL.F₂Cube, FABL.IsSplitCharacter, FABL.NonadaptiveFunctionTester.acceptanceProbability, FABL.NonadaptiveFunctionTester.transcript, FABL.NonadaptiveFunctionTester.transcript_apply, FABL.Sign, FABL.f₂DotProduct, FABL.pmfExpectation, FABL.pmfExpectation_const_one, FABL.signEncode_add, FABL.splitCharacterFunctionTester, FABL.splitCharacterTranscriptAccept, FABL.𝔽₂}
\difficulty{5}
Let $H$ be a subspace of the binary cube $\bf_2^n$, and let a split identify that cube
linearly with $H\times\bbf_2^k$, carrying $(h,0)$ to $h$. For $x,y\in\bbf_2^k$ and $h\in
H$, consider the three points corresponding to $(h,x)$, $(0,y)$, and $(0,x+y)$; the
tester samples $(x,y,h)$ uniformly and accepts exactly when the product of the first two
function values equals the third. Every function of the form
$z\mapsto(-1)^{\langle\gamma,\pi_2(\Phi^{-1}z)\rangle}$, for some $\gamma\in\bbf_2^k$,
is accepted with probability $1$.
\end{theorem}

\begin{theorem}[Sign product indicator identity]
\lean{FABL.splitCharacter_indicator_eq}
\label{FABL.splitCharacter_indicator_eq}
\leanok
\uses{FABL.Sign, FABL.signValue}
\difficulty{4}
Let $a,b,c$ be units in $$ (so each is $1$ or $-1$), and write $(a)$ for the
corresponding real value of a sign. Then the indicator of the condition $ab=c$ is
$(1+(a)(b)(c))/2$.
\end{theorem}

\begin{definition}[Split character correlation]
\lean{FABL.splitCharacterCorrelation}
\statementsource{Analysis-of-Boolean-Functions-by-Ryan-ODonnell}{pdf-ffbe9c7b3306-p197-b009}
\label{FABL.splitCharacterCorrelation}
\leanok
\uses{FABL.BinarySubspaceSplit, FABL.BinarySubspaceSplit.point, FABL.F₂Cube, FABL.Sign, FABL.signValue, FABL.𝔽₂}
Let $n$ be a natural number, let $H$ be a linear subspace of $(br/	ext{?})^n$ over the
two-element field, and let a binary subspace split of $H$ have complementary dimension
$d$. For a function $f$ from the $n$-dimensional binary cube to the signs $\{\pm1\}$,
its split character correlation is the average over independent uniform $a,b$ in the
$d$-dimensional binary cube and uniform $h\in H$ of the product of the sign values of
$f$ at the three points corresponding to $(h,a)$, $(0,b)$, and $(0,a+b)$ under the split
coordinates.
\end{definition}

\begin{theorem}[Average of an affine function]
\lean{FABL.expect_one_add_div_two}
\label{FABL.expect_one_add_div_two}
\leanok
\difficulty{1}
Let $\Omega$ be a finite nonempty set and let $g:\Omega\to\mathbb R$. The uniform
average of $(1+g(x))/2$ over $x\in\Omega$ equals one half of $1$ plus the uniform
average of $g$ over $\Omega$.
\end{theorem}

\begin{theorem}[Acceptance probability of the split character tester]
\lean{FABL.splitCharacterFunctionTester_acceptanceProbability}
\label{FABL.splitCharacterFunctionTester_acceptanceProbability}
\leanok
\uses{FABL.BinarySubspaceSplit, FABL.BinarySubspaceSplit.point, FABL.F₂Cube, FABL.NonadaptiveFunctionTester.acceptanceProbability, FABL.NonadaptiveFunctionTester.transcript_apply, FABL.Sign, FABL.SplitCharacterSeed, FABL.expect_one_add_div_two, FABL.pmfExpectation, FABL.pmfExpectation_uniformPMF_eq_expect, FABL.signValue, FABL.splitCharacterCorrelation, FABL.splitCharacterFunctionTester, FABL.splitCharacterTranscriptAccept, FABL.splitCharacter_indicator_eq, FABL.uniformPMF, FABL.𝔽₂}
\difficulty{5}
Let $H$ be a subspace of $\mathbb F_2^n$ and choose a linear coordinate isomorphism
identifying $H\times\mathbb F_2^d$ with $\mathbb F_2^n$ that sends $(h,0)$ to $h$.
Consider the tester that samples $a,b$ uniformly from $\mathbb F_2^d$ and $h$ uniformly
from $H$, queries $f$ at $(h,a)$, $(0,b)$, and $(0,a+b)$ in these coordinates, and
accepts exactly when the product of the first two sign values equals the third. For
every $f:\mathbb F_2^n\to\{\pm1\}$, its acceptance probability is $(1+C)/2$, where $C$
is the expectation over independent uniform $a,b,h$ of $f(h,a)f(0,b)f(0,a+b)$.
\end{theorem}

\begin{theorem}[A Fourier coefficient bounding a convolution average]
\lean{FABL.exists_vectorFourierCoeff_ge_expect_mul_signConvolution}
\proofsource{Analysis-of-Boolean-Functions-by-Ryan-ODonnell}{pdf-ffbe9c7b3306-p032-b004}
\label{FABL.exists_vectorFourierCoeff_ge_expect_mul_signConvolution}
\leanok
\uses{FABL.F₂Cube, FABL.Sign, FABL.convolution, FABL.signValue, FABL.sum_sq_vectorFourierCoeff_signValue_eq_one, FABL.vectorFourierCoeff, FABL.vectorFourierCoeff_convolution, FABL.vector_plancherel}
\difficulty{4}
Let $u:\mathbb F_2^d\to\mathbb R$ and let $v:\mathbb F_2^d\to\{-1,1\}$. There is a
vector $\gamma\in\mathbb F_2^d$ such that the uniform average of $u(x)$ times the
normalized additive convolution of $v$ with itself at $x$ is at most the vector Fourier
coefficient of $u$ indexed by $\gamma$.
\end{theorem}

\begin{definition}[Average on a coset]
\lean{FABL.BinarySubspaceSplit.cosetAverage}
\statementsource{Analysis-of-Boolean-Functions-by-Ryan-ODonnell}{
  pdf-ffbe9c7b3306-p077-b005,
  pdf-ffbe9c7b3306-p078-b004,
  pdf-ffbe9c7b3306-p197-b009
}
\label{FABL.BinarySubspaceSplit.cosetAverage}
\leanok
\uses{FABL.BinarySubspaceSplit, FABL.BinarySubspaceSplit.point, FABL.F₂Cube, FABL.Sign, FABL.signValue, FABL.𝔽₂}
Let $H$ be a subspace of the binary cube $\mathbb{F}_2^n$, let a binary subspace split
of $H$ be given, let $f$ map the binary cube to the signs, and let $a$ be a point of the
complementary binary cube specified by the split. The coset average at $a$ is the
average, over uniformly chosen $h\in H$, of the real sign value of $f$ evaluated at the
point determined by $h$ and $a$ under the split.
\end{definition}

\begin{definition}[Section of a sign-valued function]
\lean{FABL.BinarySubspaceSplit.sectionSign}
\statementsource{Analysis-of-Boolean-Functions-by-Ryan-ODonnell}{pdf-ffbe9c7b3306-p197-b009}
\label{FABL.BinarySubspaceSplit.sectionSign}
\leanok
\uses{FABL.BinarySubspaceSplit, FABL.BinarySubspaceSplit.point, FABL.F₂Cube, FABL.Sign, FABL.𝔽₂}
Let $H$ be a subspace of the binary cube, and choose a linear coordinate decomposition
of the cube into $H$ and a complementary binary cube. For a sign-valued function $f$ on
the full cube and a point $a$ in the complementary cube, its section at $a$ is the value
of $f$ at the point with zero $H$-coordinate and complementary coordinate $a$.
\end{definition}

\begin{theorem}[Split correlation as an averaged convolution]
\lean{FABL.splitCharacterCorrelation_eq_expect_mul_convolution}
\label{FABL.splitCharacterCorrelation_eq_expect_mul_convolution}
\leanok
\uses{FABL.BinarySubspaceSplit, FABL.BinarySubspaceSplit.cosetAverage, FABL.BinarySubspaceSplit.point, FABL.BinarySubspaceSplit.sectionSign, FABL.F₂Cube, FABL.Sign, FABL.convolution, FABL.convolution_apply_add, FABL.signValue, FABL.splitCharacterCorrelation, FABL.𝔽₂}
\difficulty{4}
Let $H$ be a linear subspace of $\mathbb F_2^n$, and suppose a linear coordinate
isomorphism identifies $H\times\mathbb F_2^d$ with $\mathbb F_2^n$ and sends $(h,0)$ to
$h$. For any $f:\mathbb F_2^n\to\{\pm1\}$, its split character correlation equals
$\mathbb E_a\bigl[A(a)(g*g)(a)\bigr]$, where $A(a)=\mathbb E_{h\in H} f(h,a)$ in these
coordinates, $g(b)=f(0,b)$, and convolution is normalized additive convolution on
$\mathbb F_2^d$.
\end{theorem}

\begin{theorem}[Vector fourier coeff coset average]
\lean{FABL.BinarySubspaceSplit.vectorFourierCoeff_cosetAverage}
\statementsource{Analysis-of-Boolean-Functions-by-Ryan-ODonnell}{
  pdf-ffbe9c7b3306-p078-b003,
  pdf-ffbe9c7b3306-p078-b004
}
\label{FABL.BinarySubspaceSplit.vectorFourierCoeff_cosetAverage}
\leanok
\uses{FABL.BinarySubspaceSplit, FABL.BinarySubspaceSplit.character, FABL.BinarySubspaceSplit.cosetAverage, FABL.BinarySubspaceSplit.point, FABL.F₂Cube, FABL.Sign, FABL.signValue, FABL.signValue_signEncode_eq_binarySign, FABL.vectorFourierCoeff, FABL.vectorFourierCoeff_eq_expect, FABL.vectorWalshCharacter, FABL.vectorWalshCharacter_apply, FABL.𝔽₂}
\difficulty{5}
% dataset-pipeline: stub -- prose pending build_informal_statements.py
Statement pending regeneration.
\end{theorem}

\begin{theorem}[Sign correlation and relative Hamming distance]
\lean{FABL.expect_signValue_mul_eq_one_sub_two_relativeHammingDist}
\proofsource{Analysis-of-Boolean-Functions-by-Ryan-ODonnell}{
  pdf-ffbe9c7b3306-p026-b003,
  pdf-ffbe9c7b3306-p026-b004
}
\label{FABL.expect_signValue_mul_eq_one_sub_two_relativeHammingDist}
\leanok
\uses{FABL.Sign, FABL.binarySignEquiv, FABL.relativeHammingDist, FABL.signEncoded_mul_eq_one_sub_two_indicator_ne, FABL.signValue, FABL.uniformProbability, FABL.uniformProbability_ne_eq_relativeHammingDist}
\difficulty{4}
Let $\Omega$ be a finite nonempty set, and let $f,g:\Omega\to\{−1,1\}$. The uniform mean
of the pointwise product $f(x)g(x)$ is $1-2d_H(f,g)/|\Omega|$, where $d_H(f,g)$ is the
number of points of $\Omega$ on which $f$ and $g$ differ.
\end{theorem}

\begin{theorem}[Vector fourier coeff coset average eq one sub two distance]
\lean{FABL.BinarySubspaceSplit.vectorFourierCoeff_cosetAverage_eq_one_sub_two_distance}
\statementsource{Analysis-of-Boolean-Functions-by-Ryan-ODonnell}{
  pdf-ffbe9c7b3306-p026-b003,
  pdf-ffbe9c7b3306-p026-b004,
  pdf-ffbe9c7b3306-p078-b003,
  pdf-ffbe9c7b3306-p078-b004
}
\label{FABL.BinarySubspaceSplit.vectorFourierCoeff_cosetAverage_eq_one_sub_two_distance}
\leanok
\uses{FABL.BinarySubspaceSplit, FABL.BinarySubspaceSplit.character, FABL.BinarySubspaceSplit.cosetAverage, FABL.BinarySubspaceSplit.vectorFourierCoeff_cosetAverage, FABL.F₂Cube, FABL.Sign, FABL.expect_signValue_mul_eq_one_sub_two_relativeHammingDist, FABL.relativeHammingDist, FABL.vectorFourierCoeff, FABL.𝔽₂}
\difficulty{1}
% dataset-pipeline: stub -- prose pending build_informal_statements.py
Statement pending regeneration.
\end{theorem}

\begin{theorem}[Split character function tester is local tester]
\lean{FABL.splitCharacterFunctionTester_isLocalTester}
\statementsource{Analysis-of-Boolean-Functions-by-Ryan-ODonnell}{
  pdf-ffbe9c7b3306-p174-b002,
  pdf-ffbe9c7b3306-p197-b009
}
\label{FABL.splitCharacterFunctionTester_isLocalTester}
\leanok
\uses{FABL.BinarySubspaceSplit, FABL.BinarySubspaceSplit.character, FABL.BinarySubspaceSplit.cosetAverage, FABL.BinarySubspaceSplit.sectionSign, FABL.BinarySubspaceSplit.vectorFourierCoeff_cosetAverage_eq_one_sub_two_distance, FABL.F₂Cube, FABL.IsClose, FABL.IsLocalTester, FABL.IsPerpendicularCharacter, FABL.exists_vectorFourierCoeff_ge_expect_mul_signConvolution, FABL.isLocalTester_of_close_of_acceptanceProbability_ge, FABL.isSplitCharacter_iff_isPerpendicularCharacter, FABL.splitCharacterCorrelation, FABL.splitCharacterCorrelation_eq_expect_mul_convolution, FABL.splitCharacterFunctionTester, FABL.splitCharacterFunctionTester_acceptanceProbability, FABL.splitCharacterFunctionTester_complete, FABL.vectorFourierCoeff, FABL.𝔽₂}
\difficulty{5}
% dataset-pipeline: stub -- prose pending build_informal_statements.py
Statement pending regeneration.
\end{theorem}

\begin{theorem}[Local testing of perpendicular characters]
\lean{FABL.perpendicularCharacterFunctionTester_isLocalTester}
\label{FABL.perpendicularCharacterFunctionTester_isLocalTester}
\leanok
\uses{FABL.F₂Cube, FABL.IsLocalTester, FABL.IsPerpendicularCharacter, FABL.binarySubspaceSplit, FABL.splitCharacterFunctionTester, FABL.splitCharacterFunctionTester_isLocalTester, FABL.𝔽₂}
\difficulty{1}
Let $n$ be a nonnegative integer, let $H$ be a subspace of $bF_2^n$, and write
$H^perp$ for its orthogonal complement under the usual dot product. Consider the class
of sign-valued functions $g$ on $bF_2^n$ of the form $g(z)=(-1)^{\langle\eta,z\rangle}$
for some $beta\in H^perp$, and the three-query tester associated with a coordinate
decomposition of $bF_2^n$ having $H$ as its first factor. The tester accepts every
function in this class with probability $1$; moreover, for every $\epsilon\in[0,1]$ and
every sign-valued function $f$ whose relative Hamming distance from each function in the
class is greater than $\epsilon$, its rejection probability is greater than $\epsilon$.
\end{theorem}

\begin{definition}[Pointwise sign twist]
\lean{FABL.signTwist}
\statementsource{Analysis-of-Boolean-Functions-by-Ryan-ODonnell}{pdf-ffbe9c7b3306-p197-b009}
\label{FABL.signTwist}
\leanok
\uses{FABL.Sign}
Let $Omega$ be a type, and let $f$ and $phase$ be sign-valued functions on $Omega$.
Their sign twist is the function that assigns to each $x$ the product $f(x)phase(x)$.
\end{definition}

\begin{theorem}[Involution of sign twisting]
\lean{FABL.signTwist_involutive}
\label{FABL.signTwist_involutive}
\leanok
\uses{FABL.Sign, FABL.signTwist}
\difficulty{4}
Let $\phi,f:\Omega\to\{\pm1\}$. Multiplying $f$ pointwise by $\phi$ twice returns $f$.
\end{theorem}

\begin{theorem}[Sign Twisting Preserves Relative Distance]
\lean{FABL.relativeHammingDist_signTwist}
\label{FABL.relativeHammingDist_signTwist}
\leanok
\uses{FABL.Sign, FABL.relativeHammingDist, FABL.signTwist}
\difficulty{3}
Let $\Omega$ be a finite nonempty set, and let $\mathrm{phase},f,g:\Omega\to\{\pm1\}$.
Multiplying both $f$ and $g$ pointwise by the same phase function preserves their
relative Hamming distance.
\end{theorem}

\begin{definition}[Twisted tester transcript]
\lean{FABL.signTwistTranscript}
\statementsource{Analysis-of-Boolean-Functions-by-Ryan-ODonnell}{pdf-ffbe9c7b3306-p197-b009}
\label{FABL.signTwistTranscript}
\leanok
\uses{FABL.NonadaptiveFunctionTester, FABL.Sign}
Let a nonadaptive tester make $r$ queries to a domain $Omega$ with sign-valued answers.
Given a seed, a transcript of optional answers, and a sign-valued phase function on the
domain, form a new transcript by multiplying each present answer by the phase at its
queried point; at positions with no query or no answer, leave the transcript absent.
\end{definition}

\begin{definition}[Sign-twisted function tester]
\lean{FABL.signTwistFunctionTester}
\statementsource{Analysis-of-Boolean-Functions-by-Ryan-ODonnell}{pdf-ffbe9c7b3306-p197-b009}
\label{FABL.signTwistFunctionTester}
\leanok
\uses{FABL.NonadaptiveFunctionTester, FABL.Sign, FABL.signTwistTranscript}
Let a nonadaptive tester query a domain $Omega$ at $r$ positions and receive
sign-valued answers, and let $phase$ assign a sign to each domain point. The
sign-twisted tester has the same seed space, seed distribution, and queries, and applies
the original acceptance rule to the transcript whose present answers are multiplied by
the phase at their queried points.
\end{definition}

\begin{theorem}[Twisting a tester transcript]
\lean{FABL.signTwistTranscript_transcript}
\label{FABL.signTwistTranscript_transcript}
\leanok
\uses{FABL.NonadaptiveFunctionTester, FABL.NonadaptiveFunctionTester.transcript, FABL.NonadaptiveFunctionTester.transcript_apply, FABL.Sign, FABL.signTwist, FABL.signTwistTranscript}
\difficulty{4}
Let $T$ be a nonadaptive function tester with a finite seed space, let $f$ and $\phi$
map its query domain to the signs $\{\pm1\}$, and fix a seed. If each answered query is
multiplied by $\phi$ at its query point, the resulting transcript is exactly the
transcript obtained by querying the pointwise product $x\mapsto f(x)\phi(x)$.
\end{theorem}

\begin{theorem}[Sign-twisted tester acceptance probability]
\lean{FABL.signTwistFunctionTester_acceptanceProbability}
\label{FABL.signTwistFunctionTester_acceptanceProbability}
\leanok
\uses{FABL.NonadaptiveFunctionTester, FABL.NonadaptiveFunctionTester.acceptanceProbability, FABL.NonadaptiveFunctionTester.transcript, FABL.Sign, FABL.signTwist, FABL.signTwistFunctionTester, FABL.signTwistTranscript, FABL.signTwistTranscript_transcript}
\difficulty{3}
Let $\Omega$ be a domain, let $r$ be a natural number, and let $T$ be a nonadaptive
tester from $\Omega$ to signs making $r$ queries. For any sign-valued functions
$\mathrm{phase},f:\Omega\to\{\pm1\}$, the acceptance probability of the tester with its
transcript values multiplied by the phase at each queried point, when run on $f$, equals
the acceptance probability of $T$ when run on the pointwise product $x\mapsto
f(x)\,\mathrm{phase}(x)$.
\end{theorem}

\begin{theorem}[Local testing under a sign twist]
\lean{FABL.isLocalTester_signTwist}
\label{FABL.isLocalTester_signTwist}
\leanok
\uses{FABL.IsFar, FABL.IsLocalTester, FABL.NonadaptiveFunctionTester, FABL.NonadaptiveFunctionTester.rejectionProbability, FABL.Sign, FABL.relativeHammingDist_signTwist, FABL.signTwist, FABL.signTwistFunctionTester, FABL.signTwistFunctionTester_acceptanceProbability, FABL.signTwist_involutive}
\difficulty{5}
Let $\Omega$ be a finite nonempty set, let $T$ be a nonadaptive tester for functions
$\Omega\to\{-1,1\}$, let $P$ be a property of such functions, and let $\rho$ be a real
rate. For any fixed function $\phi:\Omega\to\{-1,1\}$, if $T$ locally tests $P$ with
rate $\rho$, then the tester that multiplies each queried answer by $\phi$ locally tests
the property consisting of functions $f$ for which the pointwise product $\phi f$
satisfies $P$, with the same rate $\rho$.
\end{theorem}

\begin{theorem}[Invariance under equivalent properties]
\lean{FABL.IsLocalTester.congr_property}
\label{FABL.IsLocalTester.congr_property}
\leanok
\uses{FABL.IsLocalTester, FABL.NonadaptiveFunctionTester}
\difficulty{3}
Let $\Omega$ be a finite nonempty domain, let $\beta$ be a finite alphabet, and let $T$
be a nonadaptive tester with finite seed space. Let $P,Q$ be properties of functions
from $\Omega$ to $\beta$ that contain exactly the same functions, and let
$\mathrm{rate}\in\mathbb R$. If $T$ is a local tester for $P$ at rate
$\mathrm{rate}$—that is, $\mathrm{rate}>0$, every function in $P$ is accepted with
probability $1$, and every function that is $\varepsilon$-far from every member of $P$
is rejected with probability greater than $\mathrm{rate}\,\varepsilon$ for each
$\varepsilon\in[0,1]$—then $T$ is a local tester for $Q$ at the same rate.
\end{theorem}

\begin{definition}[Binary cube sign character]
\lean{FABL.f₂SignCharacter}
\leanok
\uses{FABL.F₂Cube, FABL.Sign, FABL.f₂DotProduct, FABL.signEncode}
\statementsource{Analysis-of-Boolean-Functions-by-Ryan-ODonnell}{
  pdf-ffbe9c7b3306-p014-b002,
  pdf-ffbe9c7b3306-p023-b001,
  pdf-ffbe9c7b3306-p148-b004,
  pdf-ffbe9c7b3306-p197-b009
}
For a frequency vector $gamma$ in the binary cube, its sign character is the function
on binary cube vectors $z$ whose value is the sign encoding of the binary dot product of
$gamma$ and $z$.
\end{definition}

\begin{theorem}[Product of binary sign characters]
\lean{FABL.signTwist_f₂SignCharacter}
\leanok
\uses{FABL.F₂Cube, FABL.f₂DotProduct, FABL.f₂SignCharacter, FABL.signEncode, FABL.signEncode_add, FABL.signTwist}
\difficulty{3}
For any $a,\gamma\in(\mathbb F_2)^n$, define the sign character associated with
$u\in(\mathbb F_2)^n$ by $z\mapsto(-1)^{\langle u,z\rangle}$, where the dot product is
taken in $\mathbb F_2$. Multiplying the character for $\gamma$ pointwise by the
character for $a$ gives the character for $\gamma+a$.
\end{theorem}

\begin{definition}[Affine frequency character]
\lean{FABL.IsAffineFrequencyCharacter}
\statementsource{Analysis-of-Boolean-Functions-by-Ryan-ODonnell}{
  pdf-ffbe9c7b3306-p031-b002,
  pdf-ffbe9c7b3306-p038-b010,
  pdf-ffbe9c7b3306-p197-b009
}
\label{FABL.IsAffineFrequencyCharacter}
\leanok
\uses{FABL.F₂Cube, FABL.Sign, FABL.binaryAffineSubspace, FABL.f₂SignCharacter, FABL.𝔽₂}
Let $V$ be a subspace of the binary cube, let $a$ be a binary cube vector, and let $f$
be a sign-valued function on that cube. The function $f$ is an affine frequency
character associated with $V$ and $a$ if it equals the sign character of some frequency
in the affine subspace determined by $V$ and $a$.
\end{definition}

\begin{theorem}[Affine characters and perpendicular twists]
\lean{FABL.isAffineFrequencyCharacter_iff_twist_isPerpendicularCharacter}
\label{FABL.isAffineFrequencyCharacter_iff_twist_isPerpendicularCharacter}
\leanok
\uses{FABL.F₂Cube, FABL.IsAffineFrequencyCharacter, FABL.IsPerpendicularCharacter, FABL.Sign, FABL.f₂SignCharacter, FABL.mem_binaryAffineSubspace_iff_add_mem, FABL.perpendicularSubspace, FABL.perpendicularSubspace_perpendicularSubspace, FABL.signTwist, FABL.signTwist_f₂SignCharacter, FABL.signTwist_involutive, FABL.𝔽₂}
\difficulty{5}
Let $V$ be a linear subspace of the $n$-dimensional vector space over $\mathbb F_2$, let
$a$ be a vector in that space, and let $f$ be a sign-valued function on the space. Then
$f$ is a character indexed by some vector in the affine subspace $a+V$ if and only if
the function obtained by multiplying $f$ pointwise by the character indexed by $a$ is a
character indexed by some vector in $(V^\perp)^\perp$, where perpendicularity is with
respect to the standard dot product over $\mathbb F_2$.
\end{theorem}

\begin{definition}[Affine character tester]
\lean{FABL.affineCharacterFunctionTester}
\statementsource{Analysis-of-Boolean-Functions-by-Ryan-ODonnell}{pdf-ffbe9c7b3306-p197-b009}
\label{FABL.affineCharacterFunctionTester}
\leanok
\uses{FABL.F₂Cube, FABL.NonadaptiveFunctionTester, FABL.Sign, FABL.binarySubspaceSplit, FABL.f₂SignCharacter, FABL.perpendicularSubspace, FABL.signTwistFunctionTester, FABL.splitCharacterFunctionTester, FABL.𝔽₂}
Given a subspace $V$ and a vector $a$ in the binary cube of dimension $n$, the affine
character tester is the three-query sign-valued tester obtained by taking the
split-character tester for the perpendicular subspace to $V$ and twisting its answers by
the sign character associated with $a$.
\end{definition}

\begin{theorem}[Local test for affine frequency characters]
\lean{FABL.affineCharacterFunctionTester_isLocalTester}
\label{FABL.affineCharacterFunctionTester_isLocalTester}
\leanok
\uses{FABL.F₂Cube, FABL.IsAffineFrequencyCharacter, FABL.IsLocalTester, FABL.IsLocalTester.congr_property, FABL.affineCharacterFunctionTester, FABL.f₂SignCharacter, FABL.isAffineFrequencyCharacter_iff_twist_isPerpendicularCharacter, FABL.isLocalTester_signTwist, FABL.perpendicularCharacterFunctionTester_isLocalTester, FABL.perpendicularSubspace, FABL.𝔽₂}
\difficulty{4}
Let $n$ be a natural number, let $V$ be a linear subspace of $(\mathbb{F}_2)^n$, and let
$a\in(\mathbb{F}_2)^n$. The three-query affine character tester accepts every
sign-valued function that is a character with frequency in the affine subspace
$a+V^\perp$, with acceptance probability $1$. Moreover, for every $\varepsilon\in[0,1]$,
if a sign-valued function has relative Hamming distance greater than $\varepsilon$ from
every such character, then its rejection probability is strictly greater than
$\varepsilon$.
\end{theorem}

\begin{definition}[Boolean cube and binary cube equivalence]
\lean{FABL.pcppBoolCubeEquiv}
\statementsource{Analysis-of-Boolean-Functions-by-Ryan-ODonnell}{pdf-ffbe9c7b3306-p197-b010}
\label{FABL.pcppBoolCubeEquiv}
\leanok
\uses{FABL.F₂Cube, FABL.pcppBoolF₂Equiv}
For each natural number $n$, the Boolean cube of dimension $n$ is equivalent to the
binary vector cube of dimension $n$, by identifying each Boolean coordinate with its
corresponding element of the field with two elements.
\end{definition}

\begin{definition}[Indexing the binary cube]
\lean{FABL.f₂CubeIndexEquiv}
\leanok
\uses{FABL.F₂Cube, FABL.binaryCubeSignEquiv, FABL.signCubeIndexEquiv}
\statementsource{Analysis-of-Boolean-Functions-by-Ryan-ODonnell}{pdf-ffbe9c7b3306-p197-b010}
For each natural number $n$, the set of length-$n$ vectors over the two-element field is
bijective with the set of integers from $0$ to $2^n-1$.
\end{definition}

\begin{definition}[Binary proof function]
\lean{FABL.f₂CubeProofFunction}
\leanok
\uses{FABL.F₂Cube, FABL.Sign, FABL.f₂CubeIndexEquiv, FABL.pcppBoolSignEquiv}
\statementsource{Analysis-of-Boolean-Functions-by-Ryan-ODonnell}{pdf-ffbe9c7b3306-p197-b010}
For each natural number $n$ and each Boolean function on the integers from $0$ to
$2^n-1$, define a function on length-$n$ vectors over the two-element field by indexing
the input vector and converting the resulting Boolean value to its corresponding sign.
\end{definition}

\begin{definition}[Encoding a cube function as a Boolean proof]
\lean{FABL.encodeF₂CubeFunctionProof}
\leanok
\uses{FABL.F₂Cube, FABL.Sign, FABL.f₂CubeIndexEquiv, FABL.pcppBoolSignEquiv}
\statementsource{Analysis-of-Boolean-Functions-by-Ryan-ODonnell}{pdf-ffbe9c7b3306-p197-b010}
For each natural number $n$ and each function from length-$n$ vectors over the
two-element field to signs, define a Boolean function on the integers from $0$ to
$2^n-1$ by converting each index back to a cube vector, evaluating the given function,
and converting its sign value to a Boolean value.
\end{definition}

\begin{theorem}[Encoding and decoding Boolean functions]
\lean{FABL.f₂CubeProofFunction_encodeF₂CubeFunctionProof}
\leanok
\uses{FABL.F₂Cube, FABL.Sign, FABL.encodeF₂CubeFunctionProof, FABL.f₂CubeProofFunction}
\difficulty{1}
For every natural number $n$ and every function $f:(\mathbb{F}_2)^n\to\{-1,1\}$, encode
$f$ as a Boolean function on the finite index set of size $2^n$ using the canonical
enumerations of the binary cube and the sign cube. Decoding this Boolean function
recovers $f$ at every input.
\end{theorem}

\begin{definition}[Truth table lift of a tester]
\lean{FABL.f₂TruthTableLiftTester}
\leanok
\uses{FABL.F₂Cube, FABL.NonadaptiveFunctionTester, FABL.Sign, FABL.f₂CubeIndexEquiv, FABL.pcppBoolSignEquiv}
\statementsource{Analysis-of-Boolean-Functions-by-Ryan-ODonnell}{pdf-ffbe9c7b3306-p197-b010}
Given a nonadaptive tester with $r$ query positions, whose domain is the length-$d$
binary cube and whose outputs are signs, and a natural number $n$, define a tester with
the same random seed distribution and $r$ query positions. It queries the truth-table
index corresponding to each original cube query, placing each address in the truth-table
part of the domain, and accepts exactly when the original tester accepts the transcript
obtained by converting each answered Boolean value to its corresponding sign.
\end{definition}

\begin{theorem}[Acceptance probability of a lifted tester]
\lean{FABL.f₂TruthTableLiftTester_acceptanceProbability}
\leanok
\uses{FABL.F₂Cube, FABL.NonadaptiveFunctionTester, FABL.NonadaptiveFunctionTester.acceptanceProbability, FABL.NonadaptiveFunctionTester.transcript_apply, FABL.Sign, FABL.f₂CubeProofFunction, FABL.f₂TruthTableLiftTester, FABL.inputProofWord}
\difficulty{4}
Let $d,n,r$ be nonnegative integers, let $T$ be a nonadaptive tester making $r$ queries
to functions from $\mathbb F_2^d$ to the sign set $\{-1,1\}$, let
$w\colon\{1,\ldots,n\}\to\{\mathrm{false},\mathrm{true}\}$, and let
$p\colon\{1,\ldots,2^d\}\to\{\mathrm{false},\mathrm{true}\}$. Form the tester on inputs
consisting of $w$ followed by the truth table of $p$, by translating each query to its
truth-table location and translating Boolean answers to signs. Its acceptance
probability on that input equals the acceptance probability of $T$ when given the
function on $\mathbb F_2^d$ encoded by $p$.
\end{theorem}

\begin{definition}[Standard basis vector over the two-element field]
\lean{FABL.f₂StandardBasis}
\leanok
\uses{FABL.F₂Cube}
For $i\in\{0,\ldots,n-1\}$, the standard basis vector $e_i\in\mathbb{F}_2^n$ has
coordinate $1$ at $i$ and coordinate $0$ at every other index.
\end{definition}

\begin{theorem}[Character on a standard basis vector]
\lean{FABL.f₂SignCharacter_standardBasis}
\leanok
\uses{FABL.F₂Cube, FABL.f₂DotProduct, FABL.f₂SignCharacter, FABL.f₂StandardBasis, FABL.signEncode}
\difficulty{1}
Let $n$ be a nonnegative integer, let $\gamma\colon\{0,1\}^n\to\mathbb F_2$, and let
$i\in\{1,\ldots,n\}$. The sign character associated with $\gamma$, evaluated at the
$i$th standard basis vector, is $(-1)^{\gamma_i}$.
\end{theorem}

\begin{definition}[Affine consistency transcript predicate]
\lean{FABL.affineConsistencyTranscriptAccept}
\statementsource{Analysis-of-Boolean-Functions-by-Ryan-ODonnell}{
  pdf-ffbe9c7b3306-p182-b001,
  pdf-ffbe9c7b3306-p182-b002,
  pdf-ffbe9c7b3306-p197-b010
}
\label{FABL.affineConsistencyTranscriptAccept}
\leanok
\uses{FABL.pcppBoolSignEquiv}
For a transcript of three optional Boolean values, accept if and only if all three
values are present and the product of the signs corresponding to the second and third
values equals the sign corresponding to the first; reject if any value is absent.
\end{definition}

\begin{definition}[Affine consistency tester]
\lean{FABL.affineConsistencyFunctionTester}
\statementsource{Analysis-of-Boolean-Functions-by-Ryan-ODonnell}{
  pdf-ffbe9c7b3306-p182-b001,
  pdf-ffbe9c7b3306-p182-b002,
  pdf-ffbe9c7b3306-p197-b010
}
\label{FABL.affineConsistencyFunctionTester}
\leanok
\uses{FABL.F₂Cube, FABL.NonadaptiveFunctionTester, FABL.affineConsistencyTranscriptAccept, FABL.f₂CubeIndexEquiv, FABL.f₂StandardBasis, FABL.uniformPMF}
For each positive natural number $n$, define a nonadaptive tester with three queries,
using a uniformly random seed consisting of a coordinate $i$ and a length-$n$
binary-cube vector $y$. It queries the input bit at coordinate $i$, and the truth-table
entries indexed by $y$ and by $y+e_i$, where $e_i$ is the standard basis vector at
coordinate $i$. It applies the affine consistency transcript predicate to the answers.
\end{definition}

\begin{definition}[Affine coordinate correction probability]
\lean{FABL.affineCoordinateCorrectionProbability}
\statementsource{Analysis-of-Boolean-Functions-by-Ryan-ODonnell}{
  pdf-ffbe9c7b3306-p182-b001,
  pdf-ffbe9c7b3306-p182-b002,
  pdf-ffbe9c7b3306-p197-b010
}
\label{FABL.affineCoordinateCorrectionProbability}
\leanok
\uses{FABL.F₂Cube, FABL.Sign, FABL.f₂StandardBasis, FABL.localCorrection, FABL.pcppBoolSignEquiv, FABL.uniformProbability}
For a function $f$ from the length-$n$ binary cube to signs, a coordinate $i$, and a
Boolean value $b$, define the affine coordinate correction probability to be the
proportion of cube vectors $y$ for which the local correction of $f$ in direction $e_i$
at $y$ equals the sign corresponding to $b$.
\end{definition}

\begin{theorem}[Acceptance as average coordinate correction]
\lean{FABL.affineConsistencyFunctionTester_acceptanceProbability}
\proofsource{Analysis-of-Boolean-Functions-by-Ryan-ODonnell}{pdf-ffbe9c7b3306-p197-b010}
\label{FABL.affineConsistencyFunctionTester_acceptanceProbability}
\leanok
\uses{FABL.NonadaptiveFunctionTester.acceptanceProbability, FABL.NonadaptiveFunctionTester.transcript_apply, FABL.affineConsistencyFunctionTester, FABL.affineConsistencyTranscriptAccept, FABL.affineCoordinateCorrectionProbability, FABL.f₂CubeProofFunction, FABL.inputProofWord, FABL.pmfExpectation_uniformPMF_eq_expect, FABL.uniformPMF, FABL.uniformProbability}
\difficulty{4}
Let $n>0$, let $w:\{1,\ldots,n\}\to\{\mathrm{false},\mathrm{true}\}$ be a Boolean word,
and let $p$ assign a Boolean proof value to each point of $(\mathbb{F}_2)^n$. The
acceptance probability of the affine consistency tester on the word and proof equals the
average, over a uniformly chosen coordinate $i$, of the probability that the local
correction of the sign-valued function encoded by $p$ at the standard basis vector in
coordinate $i$ agrees with the sign encoded by $w(i)$.
\end{theorem}

\begin{theorem}[Sign character as support parity]
\lean{FABL.f₂SignCharacter_eq_binaryParitySign}
\leanok
\uses{FABL.F₂Cube, FABL.binaryParitySign, FABL.f₂DotProduct_eq_coordinateSum_f₂Support, FABL.f₂SignCharacter, FABL.f₂Support}
\proofsource{Analysis-of-Boolean-Functions-by-Ryan-ODonnell}{pdf-ffbe9c7b3306-p023-b001}
\difficulty{1}
Let $n$ be a nonnegative integer and let $\gamma\colon\{0,1\}^n\to\mathbb F_2$. The sign
character $z\mapsto(-1)^{\sum_{i=1}^n\gamma_i z_i}$ equals the sign function that is
$(-1)^{\sum_{i\in\operatorname{supp}(\gamma)}z_i}$, with both sums computed in $\mathbb
F_2$.
\end{theorem}

\begin{theorem}[Local correction f₂sign character]
\lean{FABL.localCorrection_f₂SignCharacter}
\leanok
\uses{FABL.F₂Cube, FABL.f₂DotProduct, FABL.f₂SignCharacter, FABL.localCorrection, FABL.signEncode_add}
\statementsource{Analysis-of-Boolean-Functions-by-Ryan-ODonnell}{
  pdf-ffbe9c7b3306-p182-b001,
  pdf-ffbe9c7b3306-p197-b010
}
\difficulty{4}
% dataset-pipeline: stub -- prose pending build_informal_statements.py
Statement pending regeneration.
\end{theorem}

\begin{theorem}[Coordinate action of the Boolean cube equivalence]
\lean{FABL.pcppBoolCubeEquiv_apply}
\label{FABL.pcppBoolCubeEquiv_apply}
\leanok
\uses{FABL.pcppBoolCubeEquiv, FABL.pcppBoolF₂Equiv}
\difficulty{1}
For every Boolean vector $w$ of length $n$ and every coordinate $i$, the $i$th
coordinate of its image in the $n$-dimensional cube over $\mathbb F_2$ is the image of
the Boolean value $w_i$ under the standard equivalence between Boolean values and
$\mathbb F_2$.
\end{theorem}

\begin{theorem}[Encoding Boolean conjunction in the binary field]
\lean{FABL.pcppBoolF₂Equiv_and}
\leanok
\uses{FABL.pcppBoolF₂Equiv}
\difficulty{4}
For any Boolean values $a$ and $b$, the image of their conjunction under the standard
equivalence between Boolean values and $\mathbb F_2$ is the product of their respective
images in $\mathbb F_2$.
\end{theorem}

\begin{theorem}[Inverse image of zero]
\lean{FABL.pcppBoolF₂Equiv_symm_zero}
\leanok
\uses{FABL.pcppBoolF₂Equiv}
\difficulty{3}
Under the equivalence between Booleans and the field with two elements, the inverse
image of $0$ is `false`.
\end{theorem}

\begin{theorem}[Inverse image of one]
\lean{FABL.pcppBoolF₂Equiv_symm_one}
\leanok
\uses{FABL.pcppBoolF₂Equiv}
\difficulty{3}
Under the equivalence between Booleans and the field with two elements, the inverse
image of $1$ is `true`.
\end{theorem}

\begin{theorem}[Sign encoding of Booleans]
\lean{FABL.pcppBoolSignEquiv_eq_signEncode}
\label{FABL.pcppBoolSignEquiv_eq_signEncode}
\leanok
\uses{FABL.pcppBoolF₂Equiv, FABL.pcppBoolSignEquiv, FABL.signEncode}
\difficulty{1}
For every Boolean $b$, its image under the Boolean-to-sign equivalence equals $(-1)^c$,
where $c$ is the image of $b$ under the Boolean-to-field equivalence.
\end{theorem}

\begin{theorem}[Affine coordinate correction probability ge]
\lean{FABL.affineCoordinateCorrectionProbability_ge}
\statementsource{Analysis-of-Boolean-Functions-by-Ryan-ODonnell}{
  pdf-ffbe9c7b3306-p182-b001,
  pdf-ffbe9c7b3306-p197-b010
}
\label{FABL.affineCoordinateCorrectionProbability_ge}
\leanok
\uses{FABL.F₂Cube, FABL.IsClose, FABL.Sign, FABL.affineCoordinateCorrectionProbability, FABL.binaryParitySign, FABL.f₂SignCharacter, FABL.f₂SignCharacter_eq_binaryParitySign, FABL.f₂SignCharacter_standardBasis, FABL.f₂StandardBasis, FABL.f₂Support, FABL.localCorrection_successProbability, FABL.pcppBoolCubeEquiv, FABL.pcppBoolCubeEquiv_apply, FABL.pcppBoolF₂Equiv, FABL.pcppBoolSignEquiv_eq_signEncode, FABL.relativeHammingDist, FABL.signEncode}
\difficulty{4}
% dataset-pipeline: stub -- prose pending build_informal_statements.py
Statement pending regeneration.
\end{theorem}

\begin{definition}[Boolean affine subspace membership]
\lean{FABL.IsBoolAffineSubspace}
\statementsource{Analysis-of-Boolean-Functions-by-Ryan-ODonnell}{
  pdf-ffbe9c7b3306-p072-b007,
  pdf-ffbe9c7b3306-p088-b003,
  pdf-ffbe9c7b3306-p197-b009
}
\label{FABL.IsBoolAffineSubspace}
\leanok
\uses{FABL.F₂Cube, FABL.binaryAffineSubspace, FABL.pcppBoolCubeEquiv, FABL.𝔽₂}
Let $V$ be a subspace of the length-$n$ binary cube, let $a$ be a length-$n$ binary
vector, and let $w$ be a Boolean vector of length $n$. The Boolean vector $w$ belongs to
the Boolean encoding of the affine subspace $a+V$.
\end{definition}

\begin{definition}[Affine subspace proof tester]
\lean{FABL.affineSubspaceProofFunctionTester}
\statementsource{Analysis-of-Boolean-Functions-by-Ryan-ODonnell}{pdf-ffbe9c7b3306-p197-b010}
\label{FABL.affineSubspaceProofFunctionTester}
\leanok
\uses{FABL.F₂Cube, FABL.NonadaptiveFunctionTester, FABL.affineCharacterFunctionTester, FABL.f₂TruthTableLiftTester, FABL.𝔽₂}
Let $n$ be a natural number, let $V$ be a linear subspace of the Boolean cube
$\mathbb{F}_2^n$, and let $a\in\mathbb{F}_2^n$. The tester on Boolean functions over
the disjoint union of the $n$ input positions and $2^n$ proof positions is the
truth-table lift of the tester associated with the affine coset $a+V$, and makes three
queries.
\end{definition}

\begin{definition}[Affine subspace PCPP system]
\lean{FABL.affineSubspacePCPPSystem}
\statementsource{Analysis-of-Boolean-Functions-by-Ryan-ODonnell}{
  pdf-ffbe9c7b3306-p179-b004,
  pdf-ffbe9c7b3306-p180-b001,
  pdf-ffbe9c7b3306-p197-b010
}
\label{FABL.affineSubspacePCPPSystem}
\leanok
\uses{FABL.F₂Cube, FABL.PCPPSystem, FABL.affineConsistencyFunctionTester, FABL.affineSubspaceProofFunctionTester, FABL.uniformMixtureFunctionTester, FABL.𝔽₂}
Let $n$ be a positive natural number, let $V$ be a linear subspace of $\mathbb{F}_2^n$,
and let $a\in\mathbb{F}_2^n$. The Boolean PCPP system with proof length $2^n$ and three
queries uses a verifier that is the uniform mixture of the affine-subspace proof tester
and the affine-consistency tester.
\end{definition}

\begin{theorem}[Acceptance probability of the affine subspace system]
\lean{FABL.affineSubspacePCPPSystem_acceptanceProbability}
\label{FABL.affineSubspacePCPPSystem_acceptanceProbability}
\leanok
\uses{FABL.F₂Cube, FABL.NonadaptiveFunctionTester.acceptanceProbability, FABL.PCPPSystem.acceptanceProbability, FABL.affineConsistencyFunctionTester, FABL.affineSubspacePCPPSystem, FABL.affineSubspaceProofFunctionTester, FABL.inputProofWord, FABL.uniformMixtureFunctionTester_acceptanceProbability, FABL.𝔽₂}
\difficulty{3}
Let $n$ be positive, let $V$ be a linear subspace of $\mathbb F_2^n$, let $a\in\mathbb
F_2^n$, and let $w\in\{0,1\}^n$ and $p\in\{0,1\}^{2^n}$. The acceptance probability of
the affine subspace PCPP system on input $w$ and proof $p$ is the average of the
acceptance probabilities of its affine-subspace proof tester and its affine-consistency
tester, each applied to the combined input and proof word.
\end{theorem}

\begin{theorem}[Completeness of the affine subspace PCPP system]
\lean{FABL.affineSubspacePCPPSystem_complete}
\proofsource{Analysis-of-Boolean-Functions-by-Ryan-ODonnell}{pdf-ffbe9c7b3306-p197-b010}
\label{FABL.affineSubspacePCPPSystem_complete}
\leanok
\uses{FABL.F₂Cube, FABL.IsBoolAffineSubspace, FABL.NonadaptiveFunctionTester.acceptanceProbability, FABL.PCPPSystem.acceptanceProbability, FABL.affineCharacterFunctionTester_isLocalTester, FABL.affineConsistencyFunctionTester, FABL.affineConsistencyFunctionTester_acceptanceProbability, FABL.affineCoordinateCorrectionProbability, FABL.affineSubspacePCPPSystem, FABL.affineSubspacePCPPSystem_acceptanceProbability, FABL.affineSubspaceProofFunctionTester, FABL.encodeF₂CubeFunctionProof, FABL.f₂CubeProofFunction_encodeF₂CubeFunctionProof, FABL.f₂SignCharacter, FABL.f₂SignCharacter_standardBasis, FABL.f₂StandardBasis, FABL.f₂TruthTableLiftTester_acceptanceProbability, FABL.inputProofWord, FABL.localCorrection, FABL.localCorrection_f₂SignCharacter, FABL.pcppBoolCubeEquiv, FABL.pcppBoolSignEquiv, FABL.pcppBoolSignEquiv_eq_signEncode, FABL.uniformProbability, FABL.𝔽₂}
\difficulty{5}
Let $n$ be positive, let $V$ be a linear subspace of $\mathbb F_2^n$, and let
$a\in\mathbb F_2^n$. If a Boolean word $w\in\{0,1\}^n$, interpreted as a vector over
$\mathbb F_2$, belongs to $a+V$, then the affine subspace PCPP system accepts $w$ with
probability $1$ when given the truth table of the sign character indexed by the vector
corresponding to $w$, encoded as a Boolean proof of length $2^n$.
\end{theorem}

\begin{theorem}[Coordinate correction probability bound]
\lean{FABL.affineCoordinateCorrectionProbability_le_one}
\label{FABL.affineCoordinateCorrectionProbability_le_one}
\leanok
\uses{FABL.F₂Cube, FABL.Sign, FABL.affineCoordinateCorrectionProbability, FABL.f₂StandardBasis, FABL.localCorrection, FABL.pcppBoolSignEquiv, FABL.uniformProbability}
\difficulty{3}
For every sign-valued function on $(\mathbb{F}_2)^n$, every coordinate $i$, and every
Boolean value $b$, the probability that the local correction in coordinate $i$ agrees
with the sign encoded by $b$, when its input is chosen uniformly from
$(\mathbb{F}_2)^n$, is at most $1$.
\end{theorem}

\begin{theorem}[Mismatch and Coordinate-Correction Rejection]
\lean{FABL.half_mismatch_le_affineCoordinateCorrection_rejection}
\proofsource{Analysis-of-Boolean-Functions-by-Ryan-ODonnell}{
  pdf-ffbe9c7b3306-p182-b002,
  pdf-ffbe9c7b3306-p182-b001
}
\label{FABL.half_mismatch_le_affineCoordinateCorrection_rejection}
\leanok
\uses{FABL.F₂Cube, FABL.Sign, FABL.affineCoordinateCorrectionProbability, FABL.affineCoordinateCorrectionProbability_le_one, FABL.f₂StandardBasis, FABL.localCorrection, FABL.uniformProbability, FABL.uniformProbability_sign_eq_one_add_neg_one_eq_one}
\difficulty{5}
Let $n$ be a natural number, let $f:\mathbb F_2^n\to\{\pm1\}$, let
$w,w_0:\{1,\ldots,n\}\to\{\mathrm{false},\mathrm{true}\}$, and choose a coordinate $i$.
Suppose that, for a uniformly random $y\in\mathbb F_2^n$, the local correction
$f(e_i+y)f(y)$ equals the sign encoded by $w_0(i)$ with probability at least $1/2$. Then
the probability that the same correction fails to equal the sign encoded by $w(i)$ is at
least $1/2$ if $w(i)\ne w_0(i)$; if they agree, the claimed lower bound is zero.
\end{theorem}

\begin{theorem}[Affine-Consistency Rejection Bound]
\lean{FABL.half_relativeHammingDist_le_affineConsistency_rejection}
\proofsource{Analysis-of-Boolean-Functions-by-Ryan-ODonnell}{pdf-ffbe9c7b3306-p197-b010}
\label{FABL.half_relativeHammingDist_le_affineConsistency_rejection}
\leanok
\uses{FABL.F₂Cube, FABL.NonadaptiveFunctionTester.acceptanceProbability, FABL.Sign, FABL.affineConsistencyFunctionTester, FABL.affineConsistencyFunctionTester_acceptanceProbability, FABL.affineCoordinateCorrectionProbability, FABL.f₂CubeProofFunction, FABL.half_mismatch_le_affineCoordinateCorrection_rejection, FABL.inputProofWord, FABL.relativeHammingDist, FABL.uniformProbability, FABL.uniformProbability_ne_eq_relativeHammingDist}
\difficulty{4}
Let $n$ be a positive natural number, let $f:\mathbb F_2^n\to\{\pm1\}$, and let
$p:\{0,1\}^n\to\{\mathrm{false},\mathrm{true}\}$ be a proof table whose induced
sign-valued function on $\mathbb F_2^n$ equals $f$. Let $w,w_0$ be Boolean strings of
length $n$, and suppose that for every coordinate $i$, the local correction of $f$ in
direction $e_i$ equals the sign encoded by $w_0(i)$ with probability at least $1/2$
under a uniform choice of the cube point. The affine-consistency tester samples a
uniform coordinate and cube point, queries the corresponding input bit and the proof at
that point and its translate by $e_i$, and accepts when the three encoded signs satisfy
the affine consistency equation. Its rejection probability on the input formed from $w$
and $p$ is at least one half the relative Hamming distance between $w$ and $w_0$.
\end{theorem}

\begin{theorem}[Distance preserved by binary encoding]
\lean{FABL.relativeHammingDist_pcppBoolCubeEquiv}
\label{FABL.relativeHammingDist_pcppBoolCubeEquiv}
\leanok
\uses{FABL.pcppBoolCubeEquiv, FABL.pcppBoolF₂Equiv, FABL.relativeHammingDist}
\difficulty{3}
Let $n$ be a positive integer, and let $w,w_0$ be Boolean strings of length $n$. Encode
each bit as the corresponding element of $\mathbb F_2$ coordinate by coordinate. The
relative Hamming distance between the encoded strings equals the relative Hamming
distance between the original strings.
\end{theorem}

\begin{theorem}[Affine subspace pcpp system close of acceptance probability ge]
\lean{FABL.affineSubspacePCPPSystem_close_of_acceptanceProbability_ge}
\statementsource{Analysis-of-Boolean-Functions-by-Ryan-ODonnell}{
  pdf-ffbe9c7b3306-p197-b010,
  pdf-ffbe9c7b3306-p180-b001
}
\label{FABL.affineSubspacePCPPSystem_close_of_acceptanceProbability_ge}
\leanok
\uses{FABL.F₂Cube, FABL.IsBoolAffineSubspace, FABL.IsClose, FABL.IsCloseToProperty, FABL.IsLocalTester.close_of_acceptanceProbability_ge, FABL.NonadaptiveFunctionTester.acceptanceProbability, FABL.PCPPSystem.acceptanceProbability, FABL.affineCharacterFunctionTester, FABL.affineCharacterFunctionTester_isLocalTester, FABL.affineConsistencyFunctionTester, FABL.affineCoordinateCorrectionProbability, FABL.affineCoordinateCorrectionProbability_ge, FABL.affineSubspacePCPPSystem, FABL.affineSubspaceProofFunctionTester, FABL.binaryAffineSubspace, FABL.component_acceptanceProbability_ge_of_uniformMixture_ge, FABL.f₂CubeProofFunction, FABL.f₂SignCharacter, FABL.f₂TruthTableLiftTester_acceptanceProbability, FABL.half_relativeHammingDist_le_affineConsistency_rejection, FABL.inputProofWord, FABL.pcppBoolCubeEquiv, FABL.relativeHammingDist, FABL.uniformMixtureFunctionTester, FABL.𝔽₂}
\difficulty{6}
% dataset-pipeline: stub -- prose pending build_informal_statements.py
Statement pending regeneration.
\end{theorem}

\begin{theorem}[Affine subspace pcpp system is pcpp]
\lean{FABL.affineSubspacePCPPSystem_isPCPP}
\statementsource{Analysis-of-Boolean-Functions-by-Ryan-ODonnell}{pdf-ffbe9c7b3306-p197-b010}
\label{FABL.affineSubspacePCPPSystem_isPCPP}
\leanok
\uses{FABL.F₂Cube, FABL.IsBoolAffineSubspace, FABL.IsPCPP, FABL.affineSubspacePCPPSystem, FABL.affineSubspacePCPPSystem_close_of_acceptanceProbability_ge, FABL.affineSubspacePCPPSystem_complete, FABL.encodeF₂CubeFunctionProof, FABL.f₂SignCharacter, FABL.isPCPP_of_exists_close_of_acceptanceProbability_ge, FABL.pcppBoolCubeEquiv, FABL.𝔽₂}
\difficulty{4}
% dataset-pipeline: stub -- prose pending build_informal_statements.py
Statement pending regeneration.
\end{theorem}

\begin{definition}[Affine subspace PCPP description]
\lean{FABL.affineSubspacePCPPDescription}
\statementsource{Analysis-of-Boolean-Functions-by-Ryan-ODonnell}{pdf-ffbe9c7b3306-p197-b010}
\label{FABL.affineSubspacePCPPDescription}
\leanok
\uses{FABL.F₂Cube, FABL.IsBoolAffineSubspace, FABL.PCPPDescription, FABL.affineSubspacePCPPSystem, FABL.affineSubspacePCPPSystem_isPCPP, FABL.𝔽₂}
Let $n$ be positive, let $V$ be a linear subspace of $\mathbb{F}_2^n$, and let
$a\in\mathbb{F}_2^n$. The affine subspace $a+V$ of the Boolean cube has a PCPP
description with proof length $2^n$, query count three, and rejection rate $1/8$. The
described system satisfies perfect completeness and rejects every proof for an input at
relative Hamming distance greater than $\varepsilon$ from the property with probability
greater than $(1/8)\varepsilon$, for every $\varepsilon\in[0,1]$.
\end{definition}

\begin{theorem}[Affine subspace membership in dimension zero]
\lean{FABL.isBoolAffineSubspace_zero}
\label{FABL.isBoolAffineSubspace_zero}
\leanok
\uses{FABL.F₂Cube, FABL.IsBoolAffineSubspace, FABL.mem_binaryAffineSubspace_iff_add_mem, FABL.pcppBoolCubeEquiv, FABL.𝔽₂}
\difficulty{4}
Let $V$ be a linear subspace of the zero-dimensional binary vector space, let $a$ be a
vector in that space, and let $w$ be a Boolean assignment on the empty coordinate set.
The binary vector corresponding to $w$ lies in the affine subspace obtained by
translating $V$ by $a$.
\end{theorem}

\begin{theorem}[Affine subspace PCPP existence]
\lean{FABL.affineSubspacePCPPDescription_or_zero}
\proofsource{Analysis-of-Boolean-Functions-by-Ryan-ODonnell}{pdf-ffbe9c7b3306-p197-b010}
\label{FABL.affineSubspacePCPPDescription_or_zero}
\leanok
\uses{FABL.F₂Cube, FABL.IsBoolAffineSubspace, FABL.PCPPDescription, FABL.affineSubspacePCPPDescription, FABL.isBoolAffineSubspace_zero, FABL.𝔽₂}
\difficulty{4}
Let $n$ be a nonnegative integer, let $V$ be a linear subspace of $\mathbb F_2^n$, and
let $a\in\mathbb F_2^n$. Either $n=0$ and every Boolean word of length $n$, viewed as a
vector over $\mathbb F_2$, belongs to the affine subspace $a+V$, or $n$ is positive and
there exists a PCPP description for the property that a Boolean word, viewed as a vector
over $\mathbb F_2$, belongs to $a+V$.
\end{theorem}

\begin{definition}[Boolean inner product equals one]
\lean{FABL.IsInnerProductModTwoOne}
\statementsource{Analysis-of-Boolean-Functions-by-Ryan-ODonnell}{
  pdf-ffbe9c7b3306-p034-b002,
  pdf-ffbe9c7b3306-p149-b003,
  pdf-ffbe9c7b3306-p151-b006,
  pdf-ffbe9c7b3306-p197-b011
}
\label{FABL.IsInnerProductModTwoOne}
\leanok
\uses{FABL.innerProductModTwoBit, FABL.pcppBoolCubeEquiv}
Let $m$ be a natural number and let $w$ be a Boolean vector of length $2m$. The inner
product modulo two of the two length-$m$ blocks of $w$ is equal to $1$ in
$\mathbb{F}_2$.
\end{definition}

\begin{definition}[Paired coordinate product]
\lean{FABL.innerProductPairBit}
\statementsource{Analysis-of-Boolean-Functions-by-Ryan-ODonnell}{pdf-ffbe9c7b3306-p197-b011}
\label{FABL.innerProductPairBit}
\leanok
\uses{FABL.pcppBoolF₂Equiv, FABL.𝔽₂}
Let $m$ be a natural number, let $w$ be a Boolean vector of length $2m$, and let $i$ be
an index with $0\le i<m$. The paired bit at $i$ is the product in $\mathbb{F}_2$ of the
Boolean values in positions $i$ and $i+m$, after identifying Boolean values with
elements of $\mathbb{F}_2$.
\end{definition}

\begin{definition}[Prefix sum of paired products]
\lean{FABL.innerProductPrefixBit}
\statementsource{Analysis-of-Boolean-Functions-by-Ryan-ODonnell}{pdf-ffbe9c7b3306-p197-b011}
\label{FABL.innerProductPrefixBit}
\leanok
\uses{FABL.innerProductPairBit, FABL.𝔽₂}
Let $m$ be a natural number, let $w$ be a Boolean vector of length $2m$, and let $k$ be
a natural number. The prefix bit is the sum in $\mathbb{F}_2$ of the paired coordinate
products at indices $i<k$ that also satisfy $i<m$.
\end{definition}

\begin{theorem}[Successive inner product prefixes]
\lean{FABL.innerProductPrefixBit_succ}
\label{FABL.innerProductPrefixBit_succ}
\leanok
\uses{FABL.innerProductPairBit, FABL.innerProductPrefixBit}
\difficulty{3}
Let $m$ be a natural number, let $w$ be a Boolean vector of length $2m$, and let $k<m$.
The prefix sum through the first $k+1$ coordinate pairs equals the prefix sum through
the first $k$ pairs plus the product in $\mathbb{F}_2$ of the two bits in pair $k$.
\end{theorem}

\begin{theorem}[Full inner product prefix]
\lean{FABL.innerProductPrefixBit_full}
\label{FABL.innerProductPrefixBit_full}
\leanok
\uses{FABL.f₂CubeBlockEquiv, FABL.f₂DotProduct, FABL.innerProductModTwoBit, FABL.innerProductPairBit, FABL.innerProductPrefixBit, FABL.pcppBoolCubeEquiv, FABL.pcppBoolCubeEquiv_apply}
\difficulty{3}
For every natural number $m$ and Boolean vector $w$ of length $2m$, the sum in
$\mathbb{F}_2$ of the products of the two bits in each of the $m$ corresponding
coordinate pairs equals the inner product modulo two of the two length-$m$ blocks
obtained by splitting $w$ into halves.
\end{theorem}

\begin{definition}[Product proof index]
\lean{FABL.innerProductProductProofIndex}
\statementsource{Analysis-of-Boolean-Functions-by-Ryan-ODonnell}{pdf-ffbe9c7b3306-p197-b011}
\label{FABL.innerProductProductProofIndex}
\leanok
Let $m$ be a natural number and let $i$ be an index with $0\le i<m$. The product proof
index is the image of $i$ in the first block of the disjoint union of an $m$-element set
and an $(m+1)$-element set, identified with an index from $0$ to $2m$.
\end{definition}

\begin{definition}[Prefix proof index]
\lean{FABL.innerProductPrefixProofIndex}
\statementsource{Analysis-of-Boolean-Functions-by-Ryan-ODonnell}{pdf-ffbe9c7b3306-p197-b011}
\label{FABL.innerProductPrefixProofIndex}
\leanok
Let $m$ be a natural number and let $k$ be an index with $0\le k\le m$. The prefix proof
index is the image of $k$ in the second block of the disjoint union of an $m$-element
set and an $(m+1)$-element set, identified with an index from $0$ to $2m$.
\end{definition}

\begin{definition}[Linear proof word construction]
\lean{FABL.innerProductLinearProof}
\statementsource{Analysis-of-Boolean-Functions-by-Ryan-ODonnell}{pdf-ffbe9c7b3306-p197-b011}
\label{FABL.innerProductLinearProof}
\leanok
\uses{FABL.innerProductPairBit, FABL.innerProductPrefixBit, FABL.pcppBoolF₂Equiv}
For a Boolean word $w$ of length $2m$, define a Boolean word of length $2m+1$ by
splitting each output position according to the decomposition into a block of length $m$
and a block of length $m+1$. At a position in the first block, use the Boolean value
associated with the corresponding pair-bit of $w$; at a position in the second block,
use the Boolean value associated with the corresponding prefix-bit of $w$.
\end{definition}

\begin{definition}[Linear check seed space]
\lean{FABL.InnerProductLinearCheckSeed}
\statementsource{Analysis-of-Boolean-Functions-by-Ryan-ODonnell}{pdf-ffbe9c7b3306-p197-b011}
\label{FABL.InnerProductLinearCheckSeed}
\leanok
For a natural number $m$, the seed space consists of either a Boolean value or a pair
$(i,b)$, where $i$ is an index in $\{0,\ldots,m-1\}$ and $b$ is a Boolean value.
\end{definition}

\begin{definition}[Linear proof check conditions]
\lean{FABL.InnerProductLinearCheckHolds}
\statementsource{Analysis-of-Boolean-Functions-by-Ryan-ODonnell}{pdf-ffbe9c7b3306-p197-b011}
\label{FABL.InnerProductLinearCheckHolds}
\leanok
\uses{FABL.InnerProductLinearCheckSeed, FABL.innerProductPrefixProofIndex, FABL.innerProductProductProofIndex, FABL.pcppBoolF₂Equiv}
Let $m$ be a natural number, let $w$ be a Boolean word of length $2m$, let $p$ be a
Boolean word of length $2m+1$, and choose a seed from the linear check seed space. If
the seed is a Boolean value $b$, the designated final-prefix position for $b$ must
contain $b$. If the seed is $(i,	ext{false})$, the designated product position must
contain the product of the two input bits at index $i$. If the seed is $(i,	ext{true})$,
the values at the designated consecutive prefix and product positions, interpreted in
the field with two elements, must satisfy $p_i+q_i=p_{i+1}$, where $p_i$ and $p_{i+1}$
denote the designated prefix values and $q_i$ the designated product value.
\end{definition}

\begin{definition}[Linear transcript acceptance rule]
\lean{FABL.innerProductLinearTranscriptAccept}
\statementsource{Analysis-of-Boolean-Functions-by-Ryan-ODonnell}{pdf-ffbe9c7b3306-p197-b011}
\label{FABL.innerProductLinearTranscriptAccept}
\leanok
\uses{FABL.InnerProductLinearCheckSeed, FABL.pcppBoolF₂Equiv}
Given a seed and three possibly missing Boolean transcript values, accept according to
the seed type. For a Boolean seed, accept exactly when the first value is present and
equals that Boolean. For a seed $(i,	ext{false})$, accept exactly when all three values
are present and the third is the Boolean product of the first two. For a seed
$(i,	ext{true})$, accept exactly when all three are present and, after interpreting them
in the field with two elements, the first plus the second equals the third.
\end{definition}

\begin{definition}[Linear PCPP verifier system]
\lean{FABL.innerProductLinearPCPPSystem}
\statementsource{Analysis-of-Boolean-Functions-by-Ryan-ODonnell}{
  pdf-ffbe9c7b3306-p179-b004,
  pdf-ffbe9c7b3306-p180-b001,
  pdf-ffbe9c7b3306-p197-b011
}
\label{FABL.innerProductLinearPCPPSystem}
\leanok
\uses{FABL.InnerProductLinearCheckSeed, FABL.PCPPSystem, FABL.innerProductLinearTranscriptAccept, FABL.innerProductPrefixProofIndex, FABL.innerProductProductProofIndex, FABL.uniformPMF}
For each positive natural number $m$, there is a Boolean proof-carrying property-testing
system whose input has length $2m$, whose proof has length $2m+1$, and whose verifier
makes three nonadaptive query slots. Its seed is uniform over the linear check seed
space. For a Boolean seed, it queries the corresponding proof prefix position in the
first slot and makes no other queries. For a seed $(i,	ext{false})$, it queries the two
input positions in the two blocks at index $i$, then the corresponding proof product
position. For a seed $(i,	ext{true})$, it queries the proof prefix positions
corresponding to $i$ and $i+1$ and the proof product position at $i$. It accepts by the
linear transcript acceptance rule.
\end{definition}

\begin{theorem}[Rejection probability as a uniform failure rate]
\lean{FABL.innerProductLinearPCPPSystem_rejectionProbability}
\label{FABL.innerProductLinearPCPPSystem_rejectionProbability}
\leanok
\uses{FABL.InnerProductLinearCheckHolds, FABL.InnerProductLinearCheckSeed, FABL.NonadaptiveFunctionTester.acceptanceProbability, FABL.PCPPSystem.acceptanceProbability, FABL.PCPPSystem.rejectionProbability, FABL.innerProductLinearPCPPSystem, FABL.innerProductLinearTranscriptAccept, FABL.innerProductPrefixProofIndex, FABL.innerProductProductProofIndex, FABL.inputProofWord, FABL.pcppBoolF₂Equiv, FABL.pmfExpectation_uniformPMF_eq_expect, FABL.uniformPMF, FABL.uniformProbability}
\difficulty{5}
For every positive natural number $m$, Boolean input $w$ of length $2m$, and Boolean
proof of length $2m+1$, the rejection probability of the inner product verifier equals
the fraction of its seeds for which the corresponding check fails. Its seeds are either
a choice of one of the two endpoint checks, or a coordinate $i<m$ together with a choice
between checking the product at coordinate $i$ and checking the prefix-sum recurrence
there.
\end{theorem}

\begin{theorem}[Checks passed by the inner product proof]
\lean{FABL.innerProductLinearProof_checks}
\label{FABL.innerProductLinearProof_checks}
\leanok
\uses{FABL.InnerProductLinearCheckHolds, FABL.InnerProductLinearCheckSeed, FABL.IsInnerProductModTwoOne, FABL.innerProductLinearProof, FABL.innerProductPairBit, FABL.innerProductPrefixBit, FABL.innerProductPrefixBit_full, FABL.innerProductPrefixBit_succ, FABL.innerProductPrefixProofIndex, FABL.innerProductProductProofIndex, FABL.pcppBoolF₂Equiv}
\difficulty{4}
Let $m$ be a natural number, let $w$ be a Boolean word of length $2m$ whose inner
product modulo two is $1$, and form its proof by recording each coordinatewise product
and each cumulative sum of those products. Then every seed check passes: the initial
cumulative sum is $0$, the final one is $1$, each recorded product equals the product of
the corresponding input bits, and consecutive cumulative sums differ by that recorded
product modulo two.
\end{theorem}

\begin{theorem}[Prefix values determined by linear checks]
\lean{FABL.innerProductPrefixProof_eq_of_checks}
\label{FABL.innerProductPrefixProof_eq_of_checks}
\leanok
\uses{FABL.InnerProductLinearCheckHolds, FABL.innerProductPairBit, FABL.innerProductPrefixBit, FABL.innerProductPrefixBit_succ, FABL.innerProductPrefixProofIndex, FABL.innerProductProductProofIndex, FABL.pcppBoolF₂Equiv, FABL.pcppBoolF₂Equiv_and, FABL.pcppBoolF₂Equiv_false}
\difficulty{5}
Let $m$ be a natural number, let $w$ be a Boolean vector of length $2m$, and let $P$ be
a Boolean proof string indexed by $m+(m+1)$ positions. Suppose every linear check holds
for every seed, where the checks include the endpoint values, the products of paired
input bits, and the successive-prefix recurrence. Then, for every $k$ with $0\le k\le
m$, the bit of $P$ at prefix index $k$, interpreted in $\mathbb{F}_2$, equals the sum in
$\mathbb{F}_2$ of the products of the first $k$ paired input bits.
\end{theorem}

\begin{theorem}[Linear checks certify a mod-two inner product]
\lean{FABL.isInnerProductModTwoOne_of_linearChecks}
\label{FABL.isInnerProductModTwoOne_of_linearChecks}
\leanok
\uses{FABL.InnerProductLinearCheckHolds, FABL.IsInnerProductModTwoOne, FABL.innerProductPrefixBit_full, FABL.innerProductPrefixProof_eq_of_checks, FABL.pcppBoolF₂Equiv, FABL.pcppBoolF₂Equiv_true}
\difficulty{3}
Let $m$ be a nonnegative integer, let $w$ be a Boolean string of length $2m$, and let
$p$ be a Boolean string of length $2m+1$. Suppose the following checks hold for every
seed: the seed selecting the first prefix position requires its bit in $p$ to be false;
the seed selecting the final prefix position requires its bit to be true; for each
$i<m$, one seed requires the product bit of $p$ at position $i$ to equal the AND of the
two corresponding bits of $w$; and another requires the two successive prefix bits of
$p$, interpreted in $\mathbb F_2$, to differ by the product bit at position $i$. Then
the inner product modulo two of the two length-$m$ blocks of $w$, interpreted as a field
element, is $1$.
\end{theorem}

\begin{theorem}[Mass of a nonempty event]
\lean{FABL.one_div_card_le_uniformProbability_of_mem}
\label{FABL.one_div_card_le_uniformProbability_of_mem}
\leanok
\uses{FABL.uniformProbability}
\difficulty{4}
Let $\Omega$ be a finite nonempty set, let $P$ be a predicate on $\Omega$, and suppose
$P$ holds for at least one element of $\Omega$. Then the uniform probability of $P$ is
at least $1/|\Omega|$.
\end{theorem}

\begin{theorem}[Lower bound on check failure probability]
\lean{FABL.one_div_checkCount_le_innerProduct_failedCheckProbability}
\label{FABL.one_div_checkCount_le_innerProduct_failedCheckProbability}
\leanok
\uses{FABL.InnerProductLinearCheckHolds, FABL.InnerProductLinearCheckSeed, FABL.one_div_card_le_uniformProbability_of_mem, FABL.uniformProbability}
\difficulty{3}
Let $m$ be a natural number, let $w$ be a Boolean string of length $2m$, and let $p$ be
a Boolean string of length $2m+1$. Choose a seed uniformly from the $2m+2$ seeds
consisting of two endpoint checks and, for each $i=0,\ldots,m-1$, two transition checks.
A seed check holds when its endpoint bit has the required value, its product bit equals
the conjunction of the corresponding two bits of $w$, or its prefix bits satisfy the
stated recurrence modulo $2$, as appropriate. If at least one seed check fails, then the
probability that a uniformly chosen seed fails is at least $1/(2m+2)$.
\end{theorem}

\begin{definition}[Binary inner product neighbor]
\lean{FABL.innerProductFlipNeighborBinary}
\statementsource{Analysis-of-Boolean-Functions-by-Ryan-ODonnell}{pdf-ffbe9c7b3306-p197-b011}
\label{FABL.innerProductFlipNeighborBinary}
\leanok
\uses{FABL.F₂Cube, FABL.f₂CubeBlockEquiv, FABL.joinF₂CubeBlocks}
For a positive natural number $m$, split a vector in the $(2m)$-dimensional cube over
the field with two elements into blocks $x,y$ of length $m$. At index $0$, set the new
first-block coordinate to $x_0y_0+1$ and the new second-block coordinate to $1$; leave
all other coordinates unchanged, then concatenate the blocks.
\end{definition}

\begin{theorem}[Inner product flip under a neighbor map]
\lean{FABL.innerProductModTwoBit_flipNeighborBinary}
\label{FABL.innerProductModTwoBit_flipNeighborBinary}
\leanok
\uses{FABL.F₂Cube, FABL.f₂CubeBlockEquiv, FABL.f₂DotProduct, FABL.innerProductFlipNeighborBinary, FABL.innerProductModTwoBit, FABL.innerProductModTwoBit_joinF₂CubeBlocks}
\difficulty{5}
Let $m>0$ and split a vector $z\in\mathbb F_2^{2m}$ into blocks $x,y\in\mathbb F_2^m$.
Set $q=x_0y_0$, replace $x_0$ by $q+1$ and $y_0$ by $1$, and leave all other coordinates
unchanged; the inner product modulo two of the resulting vector is the original inner
product plus $1$ in $\mathbb F_2$.
\end{theorem}

\begin{theorem}[Distance to an inner-product neighbor]
\lean{FABL.hammingDist_innerProductFlipNeighborBinary_le_two}
\label{FABL.hammingDist_innerProductFlipNeighborBinary_le_two}
\leanok
\uses{FABL.F₂Cube, FABL.f₂CubeBlockEquiv, FABL.innerProductFlipNeighborBinary, FABL.joinF₂CubeBlocks, FABL.joinF₂CubeBlocks_castAdd, FABL.joinF₂CubeBlocks_natAdd}
\difficulty{5}
Let $m\ge 1$ and write $z=(x,y)$ with $x,y\in\mathbb F_2^m$. Let $p=x_0y_0$. Form $z'$
by replacing $x_0$ with $p+1$ and $y_0$ with $1$, leaving all other coordinates
unchanged. Then the Hamming distance from $z$ to $z'$ is at most $2$.
\end{theorem}

\begin{definition}[Boolean inner product neighbor]
\lean{FABL.innerProductFlipNeighbor}
\statementsource{Analysis-of-Boolean-Functions-by-Ryan-ODonnell}{pdf-ffbe9c7b3306-p197-b011}
\label{FABL.innerProductFlipNeighbor}
\leanok
\uses{FABL.innerProductFlipNeighborBinary, FABL.pcppBoolCubeEquiv}
For a positive natural number $m$, convert a Boolean word of length $2m$ to a vector
over the field with two elements, apply the binary inner product neighbor operation, and
convert the result back to a Boolean word of length $2m$.
\end{definition}

\begin{theorem}[Encoding a flipped binary neighbor]
\lean{FABL.pcppBoolCubeEquiv_innerProductFlipNeighbor}
\label{FABL.pcppBoolCubeEquiv_innerProductFlipNeighbor}
\leanok
\uses{FABL.innerProductFlipNeighbor, FABL.innerProductFlipNeighborBinary, FABL.pcppBoolCubeEquiv}
\difficulty{1}
Let $m$ be a positive natural number and let $w$ be a Boolean vector of length $2m$.
Encoding the Boolean vector obtained by the inner-product flip-neighbor operation gives
exactly the binary vector obtained by applying the corresponding operation over $\mathbb
F_2$ to the encoding of $w$.
\end{theorem}

\begin{theorem}[A neighboring word has inner product one]
\lean{FABL.innerProductFlipNeighbor_isInnerProductModTwoOne}
\label{FABL.innerProductFlipNeighbor_isInnerProductModTwoOne}
\leanok
\uses{FABL.IsInnerProductModTwoOne, FABL.innerProductFlipNeighbor, FABL.innerProductModTwoBit, FABL.innerProductModTwoBit_flipNeighborBinary, FABL.pcppBoolCubeEquiv, FABL.pcppBoolCubeEquiv_innerProductFlipNeighbor}
\difficulty{4}
Let $m\ge1$ and let $w$ be a Boolean word of length $2m$. Split its corresponding vector
over $\mathbb F_2$ into blocks $x,y\in\mathbb F_2^m$. Form a neighboring word by
replacing the first entries by $x_0y_0+1$ and $1$, respectively, and leaving all other
entries unchanged. If the inner product of the two original blocks is not $1$ in
$\mathbb F_2$, then the inner product of the two blocks in the neighboring word is $1$.
\end{theorem}

\begin{theorem}[Distance to the inner-product flip neighbor]
\lean{FABL.relativeHammingDist_innerProductFlipNeighbor_le_one_div}
\label{FABL.relativeHammingDist_innerProductFlipNeighbor_le_one_div}
\leanok
\uses{FABL.hammingDist_innerProductFlipNeighborBinary_le_two, FABL.innerProductFlipNeighbor, FABL.innerProductFlipNeighborBinary, FABL.pcppBoolCubeEquiv, FABL.pcppBoolCubeEquiv_innerProductFlipNeighbor, FABL.relativeHammingDist, FABL.relativeHammingDist_pcppBoolCubeEquiv}
\difficulty{3}
Let $m$ be a positive integer and let $w$ be a Boolean string of length $2m$. Split its
coordinates into two blocks, interpret the bits as elements of $\mathbb F_2$, and let
$x,y\in\mathbb F_2^m$ be those blocks. Form a new string by replacing the coordinate
$x_0$ with $x_0y_0+1$ and $y_0$ with $1$, leaving all other coordinates unchanged, and
then interpreting the resulting bits as Boolean values. The relative Hamming distance
from $w$ to this new string is at most $1/m$.
\end{theorem}

\begin{theorem}[Inner product PCPP guarantee]
\lean{FABL.innerProductLinearPCPPSystem_isPCPP}
\label{FABL.innerProductLinearPCPPSystem_isPCPP}
\leanok
\uses{FABL.InnerProductLinearCheckHolds, FABL.IsFar, FABL.IsInnerProductModTwoOne, FABL.IsPCPP, FABL.PCPPSystem.rejectionProbability, FABL.innerProductFlipNeighbor, FABL.innerProductFlipNeighbor_isInnerProductModTwoOne, FABL.innerProductLinearPCPPSystem, FABL.innerProductLinearPCPPSystem_rejectionProbability, FABL.innerProductLinearProof, FABL.innerProductLinearProof_checks, FABL.isInnerProductModTwoOne_of_linearChecks, FABL.one_div_checkCount_le_innerProduct_failedCheckProbability, FABL.relativeHammingDist, FABL.relativeHammingDist_innerProductFlipNeighbor_le_one_div, FABL.uniformProbability}
\difficulty{5}
For every positive natural number $m$, there is a PCPP verifier for Boolean words $w$ of
length $2m$ whose inner product modulo two is $1$, with rejection rate $1/4$. Every word
satisfying the property has a proof accepted with probability $1$; for every
$	ext{error parameter }\varepsilon\in[0,1]$, every word whose relative Hamming distance
from each satisfying word is greater than $\varepsilon$ is rejected, for every proof,
with probability strictly greater than $\varepsilon/4$.
\end{theorem}

\begin{definition}[Linear inner product PCPP description]
\lean{FABL.innerProductLinearPCPPDescription}
\statementsource{Analysis-of-Boolean-Functions-by-Ryan-ODonnell}{pdf-ffbe9c7b3306-p197-b011}
\label{FABL.innerProductLinearPCPPDescription}
\leanok
\uses{FABL.IsInnerProductModTwoOne, FABL.PCPPDescription, FABL.innerProductLinearPCPPSystem, FABL.innerProductLinearPCPPSystem_isPCPP}
For each positive natural number $m$, the described Boolean property tester uses proofs
of length $2m+1$, makes three queries, and has rejection rate $1/4$. It tests the
property that the inner product over the field with two elements of the two length-$m$
blocks of the input equals $1$.
\end{definition}

\begin{theorem}[Proof length bound]
\lean{FABL.innerProductLinearPCPPDescription_proofLength_le_twice_inputLength}
\label{FABL.innerProductLinearPCPPDescription_proofLength_le_twice_inputLength}
\leanok
\uses{FABL.innerProductLinearPCPPDescription}
\difficulty{3}
For every positive natural number $m$, the proof length in the inner product PCPP
description on inputs of length $2m$ is at most $4m$.
\end{theorem}

\begin{theorem}[The zero-dimensional inner-product bit]
\lean{FABL.not_isInnerProductModTwoOne_zero}
\label{FABL.not_isInnerProductModTwoOne_zero}
\leanok
\uses{FABL.IsInnerProductModTwoOne, FABL.f₂DotProduct, FABL.innerProductModTwoBit}
\difficulty{1}
For the unique Boolean input of length zero, split into two blocks of length zero, the
inner product of the blocks over the field with two elements is not equal to $1$.
\end{theorem}

\begin{theorem}[Existence of an inner product PCPP description]
\lean{FABL.innerProductLinearPCPPDescription_or_zero}
\label{FABL.innerProductLinearPCPPDescription_or_zero}
\leanok
\uses{FABL.IsInnerProductModTwoOne, FABL.PCPPDescription, FABL.innerProductLinearPCPPDescription, FABL.not_isInnerProductModTwoOne_zero}
\difficulty{4}
For every natural number $m$, either $m=0$ and no Boolean word of length $2m$ has inner
product modulo two equal to $1$, or $m+m$ is positive and there exists a PCPP
description for the Boolean words of length $2m$ whose inner product modulo two is $1$.
\end{theorem}

\begin{definition}[Complete quadratic value equal to one]
\lean{FABL.IsCompleteQuadraticOne}
\statementsource{Analysis-of-Boolean-Functions-by-Ryan-ODonnell}{
  pdf-ffbe9c7b3306-p034-b009,
  pdf-ffbe9c7b3306-p152-b002,
  pdf-ffbe9c7b3306-p197-b012
}
\label{FABL.IsCompleteQuadraticOne}
\leanok
\uses{FABL.completeQuadraticBit, FABL.pcppBoolCubeEquiv}
For a Boolean vector $w \colon \mathrm{Fin}(n) \to
\{\mathrm{false},\mathrm{true}\}$, say that it has complete quadratic value one
exactly when the complete quadratic bit of its image under the Boolean cube encoding
determined by $w$ is equal to $1$.
\end{definition}

\begin{definition}[Complete quadratic derivative bit]
\lean{FABL.completeQuadraticDerivativeBit}
\statementsource{Analysis-of-Boolean-Functions-by-Ryan-ODonnell}{pdf-ffbe9c7b3306-p197-b012}
\label{FABL.completeQuadraticDerivativeBit}
\leanok
\uses{FABL.F₂Cube, FABL.𝔽₂}
For a positive integer $n$, let $x\colon\{0,\ldots,n-1\}\to\mathbb{F}_2$ and choose
$i\in\{0,\ldots,n-1\}$. Define the complete quadratic derivative bit at $i$ to be the
sum in $\mathbb{F}_2$ of the coordinates $x(j)$ over all $j\ne i$.
\end{definition}

\begin{theorem}[Quadratic bit on a single-coordinate vector]
\lean{FABL.completeQuadraticBit_single}
\label{FABL.completeQuadraticBit_single}
\leanok
\uses{FABL.F₂Cube, FABL.completeQuadraticBit, FABL.𝔽₂}
\difficulty{4}
For every nonnegative integer $n$, every coordinate $i\in\{0,\ldots,n-1\}$, and every
$c\in\mathbb{F}_2$, the complete quadratic bit of the vector that is $c$ at coordinate
$i$ and zero elsewhere is $0$.
\end{theorem}

\begin{theorem}[Polar form on a single-coordinate vector]
\lean{FABL.completeQuadraticPolar_single_right}
\label{FABL.completeQuadraticPolar_single_right}
\leanok
\uses{FABL.F₂Cube, FABL.completeQuadraticDerivativeBit, FABL.completeQuadraticPolar, FABL.completeQuadraticPolarFrequency, FABL.completeQuadraticPolar_eq_dotProduct, FABL.f₂DotProduct, FABL.𝔽₂}
\difficulty{5}
Let $n$ be a natural number, $x\in\mathbb F_2^n$, $i<n$, and $c\in\mathbb F_2$. The
polar form of $x$ and the vector supported only at coordinate $i$ with value $c$ is $c$
times the sum of the coordinates $x_j$ for $j\ne i$.
\end{theorem}

\begin{theorem}[Single-coordinate addition law]
\lean{FABL.completeQuadraticBit_add_single}
\label{FABL.completeQuadraticBit_add_single}
\leanok
\uses{FABL.F₂Cube, FABL.completeQuadraticBit, FABL.completeQuadraticBit_add, FABL.completeQuadraticBit_single, FABL.completeQuadraticDerivativeBit, FABL.completeQuadraticPolar_single_right, FABL.𝔽₂}
\difficulty{3}
Let $n$ be a nonnegative integer, let $x\in\mathbb{F}_2^n$, let $i$ be a coordinate, and
let $c\in\mathbb{F}_2$. Adding $c$ to coordinate $i$ changes the complete quadratic bit
by $c\sum_{j\ne i}x_j$; that is, the new value is the original value plus this quantity
in $\mathbb{F}_2$.
\end{theorem}

\begin{theorem}[Derivative under a coordinate change]
\lean{FABL.completeQuadraticDerivativeBit_add_single_of_ne}
\label{FABL.completeQuadraticDerivativeBit_add_single_of_ne}
\leanok
\uses{FABL.F₂Cube, FABL.completeQuadraticDerivativeBit, FABL.𝔽₂}
\difficulty{4}
Let $x\in\mathbb{F}_2^n$, let $i,j\in\{0,\ldots,n-1\}$ be distinct, and let
$c\in\mathbb{F}_2$. The derivative bit at coordinate $j$, formed by summing all
coordinates of its input other than $j$, changes by adding $c$ when $c$ is added to
coordinate $i$ of $x$.
\end{theorem}

\begin{theorem}[Two-coordinate addition law]
\lean{FABL.completeQuadraticBit_add_two_single}
\label{FABL.completeQuadraticBit_add_two_single}
\leanok
\uses{FABL.F₂Cube, FABL.completeQuadraticBit, FABL.completeQuadraticBit_add_single, FABL.completeQuadraticDerivativeBit, FABL.completeQuadraticDerivativeBit_add_single_of_ne}
\difficulty{3}
Let $n$ be a nonnegative integer and let $x\in\mathbb{F}_2^n$. For distinct coordinates
$i,j$, adding $1$ to both coordinates changes the complete quadratic bit to its original
value plus $\sum_{k\ne i}x_k+\sum_{k\ne j}x_k+1$, with the sum taken in $\mathbb{F}_2$.
\end{theorem}

\begin{theorem}[Only two elements of the binary field]
\lean{FABL.f₂_eq_zero_of_ne_one}
\leanok
\uses{FABL.𝔽₂}
\difficulty{4}
If $x\in\mathbb F_2$ and $x\ne1$, then $x=0$.
\end{theorem}

\begin{definition}[Quadratic flip neighbor]
\lean{FABL.completeQuadraticFlipNeighborBinary}
\label{FABL.completeQuadraticFlipNeighborBinary}
\leanok
\uses{FABL.F₂Cube, FABL.completeQuadraticBit, FABL.completeQuadraticDerivativeBit}
For $n\ge 2$ and $x\in\mathbb F_2^n$, define a neighboring cube point as follows: if the
complete quadratic bit of $x$ is $1$, return $x$. Otherwise, if some coordinate has
complete quadratic derivative bit $1$, flip one such coordinate; if none does, flip
coordinates $0$ and $1$.
\end{definition}

\begin{theorem}[Quadratic value after a coordinate repair]
\lean{FABL.completeQuadraticFlipNeighborBinary_is_one}
\label{FABL.completeQuadraticFlipNeighborBinary_is_one}
\leanok
\uses{FABL.F₂Cube, FABL.completeQuadraticBit, FABL.completeQuadraticBit_add_single, FABL.completeQuadraticBit_add_two_single, FABL.completeQuadraticDerivativeBit, FABL.completeQuadraticFlipNeighborBinary, FABL.f₂_eq_zero_of_ne_one}
\difficulty{5}
Let $n\ge 2$ and let $x\in\mathbb{F}_2^n$. Form a vector $y$ as follows: if the complete
quadratic bit of $x$ is $1$, set $y=x$; otherwise, if some coordinate has derivative bit
$1$, flip one such coordinate; if no such coordinate exists, flip coordinates $0$ and
$1$. Then the complete quadratic bit of $y$ is $1$.
\end{theorem}

\begin{theorem}[Distance after one coordinate change]
\lean{FABL.hammingDist_add_single_le_one}
\label{FABL.hammingDist_add_single_le_one}
\leanok
\uses{FABL.F₂Cube, FABL.𝔽₂}
\difficulty{4}
For every $n\ge 0$, every vector $x\in\mathbb F_2^n$, every coordinate $i$, and every
$c\in\mathbb F_2$, the Hamming distance from $x$ to the vector obtained by adding $c$ at
coordinate $i$ is at most $1$.
\end{theorem}

\begin{theorem}[Distance after two coordinate changes]
\lean{FABL.hammingDist_add_two_single_le_two}
\label{FABL.hammingDist_add_two_single_le_two}
\leanok
\uses{FABL.F₂Cube}
\difficulty{4}
For every $n\ge 0$, every $x\in\mathbb F_2^n$, and every pair of coordinates $i,j$, the
Hamming distance from $x$ to the vector obtained by adding $1$ at both coordinates is at
most $2$.
\end{theorem}

\begin{theorem}[Distance to a quadratic flip neighbor]
\lean{FABL.hammingDist_completeQuadraticFlipNeighborBinary_le_two}
\label{FABL.hammingDist_completeQuadraticFlipNeighborBinary_le_two}
\leanok
\uses{FABL.F₂Cube, FABL.completeQuadraticFlipNeighborBinary, FABL.hammingDist_add_single_le_one, FABL.hammingDist_add_two_single_le_two}
\difficulty{4}
Let $n\ge 2$ and let $x\in\mathbb F_2^n$. The Hamming distance between $x$ and its
complete quadratic flip neighbor is at most $2$.
\end{theorem}

\begin{definition}[Complete quadratic flip neighbor]
\lean{FABL.completeQuadraticFlipNeighbor}
\statementsource{Analysis-of-Boolean-Functions-by-Ryan-ODonnell}{pdf-ffbe9c7b3306-p197-b012}
\label{FABL.completeQuadraticFlipNeighbor}
\leanok
\uses{FABL.completeQuadraticFlipNeighborBinary, FABL.pcppBoolCubeEquiv}
For an integer $n\ge 2$, let $E_n$ denote the Boolean-cube to binary-cube equivalence,
and let $Q_n$ denote the complete quadratic flip-neighbor map on the binary cube. The
complete quadratic flip neighbor of $w\in\{0,1\}^n$ is $E_n^{-1}(Q_n(E_n(w)))$.
\end{definition}

\begin{theorem}[Boolean repair satisfies the quadratic condition]
\lean{FABL.completeQuadraticFlipNeighbor_isCompleteQuadraticOne}
\label{FABL.completeQuadraticFlipNeighbor_isCompleteQuadraticOne}
\leanok
\uses{FABL.IsCompleteQuadraticOne, FABL.completeQuadraticFlipNeighbor, FABL.completeQuadraticFlipNeighborBinary_is_one, FABL.pcppBoolCubeEquiv}
\difficulty{3}
Let $n\ge 2$ and let $w$ be a Boolean vector of length $n$. Convert $w$ coordinatewise
to a vector in $\mathbb{F}_2^n$, apply the repair that leaves a vector with complete
quadratic bit $1$ unchanged, otherwise flips a coordinate with derivative bit $1$ if one
exists, and otherwise flips coordinates $0$ and $1$, and convert the result back to a
Boolean vector $w'$. The complete quadratic bit of the encoded vector $w'$ is $1$.
\end{theorem}

\begin{theorem}[Distance to a quadratic flip neighbor]
\lean{FABL.relativeHammingDist_completeQuadraticFlipNeighbor_le_two_div}
\label{FABL.relativeHammingDist_completeQuadraticFlipNeighbor_le_two_div}
\leanok
\uses{FABL.completeQuadraticFlipNeighbor, FABL.completeQuadraticFlipNeighborBinary, FABL.hammingDist_completeQuadraticFlipNeighborBinary_le_two, FABL.pcppBoolCubeEquiv, FABL.relativeHammingDist, FABL.relativeHammingDist_pcppBoolCubeEquiv}
\difficulty{3}
Let $n$ be a nonzero integer with $2\le n$, and let $w$ be a Boolean vector of length
$n$. The relative Hamming distance from $w$ to its complete quadratic flip neighbor is
at most $2/n$.
\end{theorem}

\begin{definition}[Complete quadratic truncation]
\lean{FABL.completeQuadraticTruncation}
\statementsource{Analysis-of-Boolean-Functions-by-Ryan-ODonnell}{
  pdf-ffbe9c7b3306-p152-b002,
  pdf-ffbe9c7b3306-p197-b012
}
\label{FABL.completeQuadraticTruncation}
\leanok
\uses{FABL.F₂Cube, FABL.pcppBoolF₂Equiv}
For a Boolean vector $w \colon \mathrm{Fin}(n) \to
\{\mathrm{false},\mathrm{true}\}$ and a natural number $k$, define its truncation
through coordinate $k$ to have coordinate $i$ equal to the encoded value of $w_i$ when
$i<k$, and equal to $0$ otherwise.
\end{definition}

\begin{definition}[Complete quadratic linear prefix bit]
\lean{FABL.completeQuadraticLinearPrefixBit}
\statementsource{Analysis-of-Boolean-Functions-by-Ryan-ODonnell}{
  pdf-ffbe9c7b3306-p152-b002,
  pdf-ffbe9c7b3306-p197-b012
}
\label{FABL.completeQuadraticLinearPrefixBit}
\leanok
\uses{FABL.pcppBoolF₂Equiv, FABL.𝔽₂}
For a Boolean vector $w \colon \mathrm{Fin}(n) \to
\{\mathrm{false},\mathrm{true}\}$ and a natural number $k$, define its linear prefix
bit to be the sum in $\mathbb{F}_2$ of the encoded values $w_i$ over all coordinates
$i<k$.
\end{definition}

\begin{definition}[Complete quadratic prefix bit]
\lean{FABL.completeQuadraticPrefixBit}
\statementsource{Analysis-of-Boolean-Functions-by-Ryan-ODonnell}{
  pdf-ffbe9c7b3306-p152-b002,
  pdf-ffbe9c7b3306-p197-b012
}
\label{FABL.completeQuadraticPrefixBit}
\leanok
\uses{FABL.completeQuadraticBit, FABL.completeQuadraticTruncation, FABL.𝔽₂}
For a Boolean vector $w \colon \mathrm{Fin}(n) \to
\{\mathrm{false},\mathrm{true}\}$ and a natural number $k$, define its complete
quadratic prefix bit to be the complete quadratic bit of the vector obtained by
truncating $w$ through coordinate $k$.
\end{definition}

\begin{theorem}[Successive quadratic truncations]
\lean{FABL.completeQuadraticTruncation_succ}
\label{FABL.completeQuadraticTruncation_succ}
\leanok
\uses{FABL.completeQuadraticTruncation, FABL.pcppBoolF₂Equiv}
\difficulty{5}
Let $n$ be a natural number, let $w:\{0,ots,n-1
ightarrow	ext{Bool}$, and let $k<n$.
Regard $w$ as an $F_2^n$-valued vector by sending each Boolean coordinate to its
corresponding element of $F_2$. The truncation retaining the first $k+1$ coordinates
equals the truncation retaining the first $k$ coordinates plus the vector supported at
coordinate $k$ with value $w(k)$.
\end{theorem}

\begin{theorem}[Prefix sum at the next coordinate]
\lean{FABL.completeQuadraticLinearPrefixBit_succ}
\label{FABL.completeQuadraticLinearPrefixBit_succ}
\leanok
\uses{FABL.completeQuadraticLinearPrefixBit, FABL.pcppBoolF₂Equiv}
\difficulty{5}
Let $w$ be a Boolean vector of length $n$ and let $k<n$. Encode each bit as an element
of $\mathbb{F}_2$, and let the prefix sum through index $k$ mean the sum of encoded
coordinates with indices strictly less than the specified bound. The prefix sum with
bound $k+1$ equals the prefix sum with bound $k$ plus the encoded bit $w_k$.
\end{theorem}

\begin{theorem}[Derivative of a truncated Boolean vector]
\lean{FABL.completeQuadraticDerivativeBit_truncation}
\label{FABL.completeQuadraticDerivativeBit_truncation}
\leanok
\uses{FABL.completeQuadraticDerivativeBit, FABL.completeQuadraticLinearPrefixBit, FABL.completeQuadraticTruncation}
\difficulty{4}
Let $w$ be a Boolean vector of length $n$, and let $k<n$. Encode its first $k$ bits as
elements of $\mathbb{F}_2$ and set all later coordinates to zero. The derivative bit of
this truncated vector at coordinate $k$ equals the sum in $\mathbb{F}_2$ of the encoded
first $k$ bits of $w$.
\end{theorem}

\begin{theorem}[Quadratic prefix recurrence]
\lean{FABL.completeQuadraticPrefixBit_succ}
\label{FABL.completeQuadraticPrefixBit_succ}
\leanok
\uses{FABL.completeQuadraticBit_add_single, FABL.completeQuadraticDerivativeBit_truncation, FABL.completeQuadraticLinearPrefixBit, FABL.completeQuadraticPrefixBit, FABL.completeQuadraticTruncation_succ, FABL.pcppBoolF₂Equiv}
\difficulty{3}
Let $n$ be a natural number, let $w$ be a Boolean sequence of length $n$, and let $k<n$.
The quadratic value of the prefix of length $k+1$ equals the quadratic value of the
prefix of length $k$ plus, in $\mathbb F_2$, the product of the linear sum of the first
$k$ mapped coordinates and the mapped value of $w_k$.
\end{theorem}

\begin{theorem}[Full quadratic prefix value]
\lean{FABL.completeQuadraticPrefixBit_full}
\label{FABL.completeQuadraticPrefixBit_full}
\leanok
\uses{FABL.completeQuadraticBit, FABL.completeQuadraticPrefixBit, FABL.completeQuadraticTruncation, FABL.pcppBoolCubeEquiv, FABL.pcppBoolCubeEquiv_apply}
\difficulty{3}
Let $n$ be a natural number and let $w$ be a Boolean sequence of length $n$. The
quadratic value of the prefix of length $n$ equals the quadratic value of the vector
obtained by mapping every Boolean coordinate of $w$ into $\mathbb F_2$.
\end{theorem}

\begin{definition}[Index in the First Summand]
\lean{FABL.completeQuadraticProductProofIndex}
\statementsource{Analysis-of-Boolean-Functions-by-Ryan-ODonnell}{pdf-ffbe9c7b3306-p197-b012}
\label{FABL.completeQuadraticProductProofIndex}
\leanok
For a natural number $n$, define a map from $\mathrm{Fin}(n)$ to
$\mathrm{Fin}(n+2(n+1))$ by sending each index to the corresponding index in the first
summand of $\mathrm{Fin}(n)\oplus\mathrm{Fin}(2(n+1))$.
\end{definition}

\begin{definition}[Index for linear-prefix proof coordinates]
\lean{FABL.completeQuadraticLinearPrefixProofIndex}
\statementsource{Analysis-of-Boolean-Functions-by-Ryan-ODonnell}{pdf-ffbe9c7b3306-p197-b012}
\label{FABL.completeQuadraticLinearPrefixProofIndex}
\leanok
For each $k\in\{0,1,\ldots,n\}$, define its linear-prefix proof index to be the position
in a sequence of length $n+2(n+1)$ obtained by placing $k$ in the first of two blocks of
length $n+1$, both of which follow an initial block of length $n$.
\end{definition}

\begin{definition}[Complete quadratic prefix proof index]
\lean{FABL.completeQuadraticPrefixProofIndex}
\statementsource{Analysis-of-Boolean-Functions-by-Ryan-ODonnell}{pdf-ffbe9c7b3306-p197-b012}
\label{FABL.completeQuadraticPrefixProofIndex}
\leanok
For $k\in\{0,\ldots,n\}$, define its index in a sequence of length $n+2(n+1)$ to be the
position $n+(n+1)+k$, with positions numbered from $0$.
\end{definition}

\begin{definition}[Quadratic linear proof word]
\lean{FABL.completeQuadraticLinearProof}
\statementsource{Analysis-of-Boolean-Functions-by-Ryan-ODonnell}{pdf-ffbe9c7b3306-p197-b012}
\label{FABL.completeQuadraticLinearProof}
\leanok
\uses{FABL.completeQuadraticLinearPrefixBit, FABL.completeQuadraticPrefixBit, FABL.pcppBoolF₂Equiv}
Given a Boolean word $w$ of length $n$, define a Boolean proof word of length $n+2(n+1)$
in three blocks. At position $i<n$, its bit is the product over $\mathbb F_2$ of the
linear-prefix bit of $w$ at $i$ and the bit $w_i$; in the next block, the bits are the
linear-prefix bits of $w$, and in the final block they are the quadratic-prefix bits of
$w$.
\end{definition}

\begin{definition}[Endpoint-check type]
\lean{FABL.CompleteQuadraticEndpointCheck}
\statementsource{Analysis-of-Boolean-Functions-by-Ryan-ODonnell}{pdf-ffbe9c7b3306-p197-b012}
\label{FABL.CompleteQuadraticEndpointCheck}
\leanok
Define the endpoint-check type to be the disjoint union of the Boolean type and a
singleton type.
\end{definition}

\begin{definition}[Transition-check type]
\lean{FABL.CompleteQuadraticTransitionCheck}
\statementsource{Analysis-of-Boolean-Functions-by-Ryan-ODonnell}{pdf-ffbe9c7b3306-p197-b012}
\label{FABL.CompleteQuadraticTransitionCheck}
\leanok
Define the transition-check type to be the disjoint union of the Boolean type and a
singleton type.
\end{definition}

\begin{definition}[Linear-check seed type]
\lean{FABL.CompleteQuadraticLinearCheckSeed}
\statementsource{Analysis-of-Boolean-Functions-by-Ryan-ODonnell}{pdf-ffbe9c7b3306-p197-b012}
\label{FABL.CompleteQuadraticLinearCheckSeed}
\leanok
\uses{FABL.CompleteQuadraticEndpointCheck, FABL.CompleteQuadraticTransitionCheck}
For a natural number $n$, define a seed type whose alternatives are either an endpoint
check, or a position $i<n$ paired with a transition check.
\end{definition}

\begin{definition}[Validity of a linear-check seed]
\lean{FABL.CompleteQuadraticLinearCheckHolds}
\statementsource{Analysis-of-Boolean-Functions-by-Ryan-ODonnell}{pdf-ffbe9c7b3306-p197-b012}
\label{FABL.CompleteQuadraticLinearCheckHolds}
\leanok
\uses{FABL.CompleteQuadraticLinearCheckSeed, FABL.completeQuadraticLinearPrefixProofIndex, FABL.completeQuadraticPrefixProofIndex, FABL.completeQuadraticProductProofIndex, FABL.pcppBoolF₂Equiv}
Let $w$ be a Boolean word of length $n$, let $P$ be a Boolean proof word of length
$n+2(n+1)$, and let a seed be either an endpoint check or a position $i<n$ paired with a
transition check. The seed holds of $P$ and $w$ precisely under the corresponding
condition: the false endpoint check requires the linear-prefix proof bit at index $0$ to
be false; the true endpoint check requires the quadratic-prefix proof bit at index $0$
to be false; the other endpoint check requires the quadratic-prefix proof bit at index
$n$ to be true. At a position $i$, the false transition check requires the linear-prefix
proof bits at $i$ and $i+1$ to sum with $w_i$ to zero in $\mathbb F_2$; the true
transition check requires the product proof bit to equal the conjunction of the
linear-prefix proof bit at $i$ and $w_i$; and the remaining transition check requires
the quadratic-prefix proof bits at $i$ and $i+1$ to sum with the product proof bit at
$i$ to zero in $\mathbb F_2$.
\end{definition}

\begin{definition}[Linear-check transcript acceptance]
\lean{FABL.completeQuadraticLinearTranscriptAccept}
\statementsource{Analysis-of-Boolean-Functions-by-Ryan-ODonnell}{pdf-ffbe9c7b3306-p197-b012}
\label{FABL.completeQuadraticLinearTranscriptAccept}
\leanok
\uses{FABL.CompleteQuadraticLinearCheckSeed, FABL.pcppBoolF₂Equiv}
Given a seed and a transcript of three optional Boolean values, define acceptance as
follows. An endpoint seed accepts exactly when the first value is present and is false
for either Boolean endpoint alternative, or true for the singleton alternative. For a
position seed with the false transition alternative, all three values must be present
and their first two Boolean encodings must sum in $\mathbb F_2$ to the third; with the
true alternative, all three must be present and the third must be the conjunction of the
first two; with the singleton alternative, all three must be present and the first and
second encodings must sum in $\mathbb F_2$ to the third. Missing required values cause
rejection.
\end{definition}

\begin{definition}[Nonadaptive quadratic linear PCPP system]
\lean{FABL.completeQuadraticLinearPCPPSystem}
\statementsource{Analysis-of-Boolean-Functions-by-Ryan-ODonnell}{pdf-ffbe9c7b3306-p197-b012}
\label{FABL.completeQuadraticLinearPCPPSystem}
\leanok
\uses{FABL.CompleteQuadraticLinearCheckSeed, FABL.PCPPSystem, FABL.completeQuadraticLinearPrefixProofIndex, FABL.completeQuadraticLinearTranscriptAccept, FABL.completeQuadraticPrefixProofIndex, FABL.completeQuadraticProductProofIndex, FABL.uniformPMF}
For each natural number $n$, define a Boolean-proof proximity proof checking system for
Boolean inputs of length $n$ and proofs of length $n+2(n+1)$, using three nonadaptive
queries. Its seed is drawn uniformly from the endpoint and position-check seeds:
endpoint seeds query the relevant single proof bit, while position seeds query the input
and proof bits needed for the selected transition check. The verifier accepts according
to the corresponding endpoint or transition condition on the returned answers; any
unused query positions are left unanswered.
\end{definition}

\begin{theorem}[Rejection probability as a uniform failure probability]
\lean{FABL.completeQuadraticLinearPCPPSystem_rejectionProbability}
\label{FABL.completeQuadraticLinearPCPPSystem_rejectionProbability}
\leanok
\uses{FABL.CompleteQuadraticLinearCheckHolds, FABL.CompleteQuadraticLinearCheckSeed, FABL.NonadaptiveFunctionTester.acceptanceProbability, FABL.PCPPSystem.acceptanceProbability, FABL.PCPPSystem.rejectionProbability, FABL.completeQuadraticLinearPCPPSystem, FABL.completeQuadraticLinearPrefixProofIndex, FABL.completeQuadraticLinearTranscriptAccept, FABL.completeQuadraticPrefixProofIndex, FABL.completeQuadraticProductProofIndex, FABL.inputProofWord, FABL.pcppBoolF₂Equiv, FABL.pmfExpectation_uniformPMF_eq_expect, FABL.uniformPMF, FABL.uniformProbability}
\difficulty{5}
Let $n\in\mathbb{N}$, let $w$ be a Boolean vector of length $n$, and let $\pi$ be a
Boolean proof word of length $n+(n+1)+(n+1)$. The rejection probability of the specified
verifier on input $w$ and proof $\pi$ equals the uniform probability, over its finite
seed space, that the corresponding check does not hold.
\end{theorem}

\begin{theorem}[Complete quadratic linear proof checks]
\lean{FABL.completeQuadraticLinearProof_checks}
\statementsource{Analysis-of-Boolean-Functions-by-Ryan-ODonnell}{pdf-ffbe9c7b3306-p197-b012}
\label{FABL.completeQuadraticLinearProof_checks}
\leanok
\uses{FABL.CompleteQuadraticLinearCheckHolds, FABL.CompleteQuadraticLinearCheckSeed, FABL.IsCompleteQuadraticOne, FABL.completeQuadraticBit, FABL.completeQuadraticLinearPrefixBit, FABL.completeQuadraticLinearPrefixBit_succ, FABL.completeQuadraticLinearPrefixProofIndex, FABL.completeQuadraticLinearProof, FABL.completeQuadraticPrefixBit, FABL.completeQuadraticPrefixBit_full, FABL.completeQuadraticPrefixBit_succ, FABL.completeQuadraticPrefixProofIndex, FABL.completeQuadraticProductProofIndex, FABL.completeQuadraticTruncation, FABL.pcppBoolF₂Equiv, FABL.pcppBoolF₂Equiv_and, FABL.pcppBoolF₂Equiv_true}
\difficulty{5}
% dataset-pipeline: stub -- prose pending build_informal_statements.py
Statement pending regeneration.
\end{theorem}

\begin{theorem}[Linear prefix values from consistency checks]
\lean{FABL.completeQuadraticLinearPrefixProof_eq_of_checks}
\label{FABL.completeQuadraticLinearPrefixProof_eq_of_checks}
\leanok
\uses{FABL.CompleteQuadraticLinearCheckHolds, FABL.completeQuadraticLinearPrefixBit, FABL.completeQuadraticLinearPrefixBit_succ, FABL.completeQuadraticLinearPrefixProofIndex, FABL.pcppBoolF₂Equiv}
\difficulty{4}
Let $n$ be a nonnegative integer, let $w_0,\ldots,w_{n-1}$ be Boolean values, and
let a Boolean table assign values to the indices for linear prefixes, quadratic
prefixes, and coordinate products. Suppose the table satisfies all of the following
checks: the linear prefix at $0$ is false; the quadratic prefix at $0$ is false and at
$n$ is true; for each $i<n$, the linear prefix at $i+1$ is the sum in $\mathbb F_2$ of
the linear prefix at $i$ and the image of $w_i$; the product entry at $i$ is the Boolean
conjunction of the linear prefix at $i$ and $w_i$; and the quadratic prefix at $i+1$ is
the sum in $\mathbb F_2$ of the quadratic prefix at $i$ and the image of that product
entry. Then for every $k\le n$, the image in $\mathbb F_2$ of the table's linear prefix
at $k$ equals the sum of the images of those $w_i$ with $i<k$.
\end{theorem}

\begin{theorem}[Quadratic prefix values from consistency checks]
\lean{FABL.completeQuadraticPrefixProof_eq_of_checks}
\label{FABL.completeQuadraticPrefixProof_eq_of_checks}
\leanok
\uses{FABL.CompleteQuadraticLinearCheckHolds, FABL.completeQuadraticBit, FABL.completeQuadraticLinearPrefixBit, FABL.completeQuadraticLinearPrefixProof_eq_of_checks, FABL.completeQuadraticPrefixBit, FABL.completeQuadraticPrefixBit_succ, FABL.completeQuadraticPrefixProofIndex, FABL.completeQuadraticProductProofIndex, FABL.completeQuadraticTruncation, FABL.pcppBoolF₂Equiv, FABL.pcppBoolF₂Equiv_and}
\difficulty{5}
Let $n$ be a natural number, let $w$ be a Boolean sequence of length $n$, and let a
Boolean table assign values to the linear prefixes, quadratic prefixes, and coordinate
products. Suppose the linear prefix at $0$ and quadratic prefix at $0$ are false, the
quadratic prefix at $n$ is true, and for every $i<n$ the linear prefix at $i+1$ is the
sum in $\mathbb F_2$ of the linear prefix at $i$ and the image of $w_i$; the product
entry at $i$ is the conjunction of the linear prefix at $i$ and $w_i$; and the quadratic
prefix at $i+1$ is the sum in $\mathbb F_2$ of the quadratic prefix at $i$ and the image
of that product entry. Then for every $k\le n$, the image in $\mathbb F_2$ of the
table's quadratic prefix at $k$ equals the quadratic value of the first $k$ mapped
coordinates of $w$.
\end{theorem}

\begin{theorem}[Quadratic parity from linear checks]
\lean{FABL.isCompleteQuadraticOne_of_linearChecks}
\label{FABL.isCompleteQuadraticOne_of_linearChecks}
\leanok
\uses{FABL.CompleteQuadraticLinearCheckHolds, FABL.IsCompleteQuadraticOne, FABL.completeQuadraticPrefixBit_full, FABL.completeQuadraticPrefixProof_eq_of_checks, FABL.pcppBoolF₂Equiv, FABL.pcppBoolF₂Equiv_true}
\difficulty{3}
Let $n$ be a natural number, let $w:	ext{Fin}(n)	o	ext{Bool}$, and let $P$ be a
Boolean-valued table indexed by $	ext{Fin}(n+2(n+1))$. Suppose that every endpoint,
linear-transition, product, and quadratic-prefix check holds for $w$ and $P$: the linear
and quadratic prefix bits start at $0$, the final quadratic prefix bit is $1$, each
linear prefix advances by adding the corresponding bit of $w$ modulo $2$, each product
bit is the conjunction of the preceding linear prefix bit and the corresponding bit of
$w$, and each quadratic prefix advances by adding the corresponding product bit modulo
$2$. Then the sum of $w_iw_j$ over $i<j$ equals $1$ in $br/2br$.
\end{theorem}

\begin{theorem}[Lower bound on check failure probability]
\lean{FABL.one_div_checkCount_le_completeQuadratic_failedCheckProbability}
\label{FABL.one_div_checkCount_le_completeQuadratic_failedCheckProbability}
\leanok
\uses{FABL.CompleteQuadraticEndpointCheck, FABL.CompleteQuadraticLinearCheckHolds, FABL.CompleteQuadraticLinearCheckSeed, FABL.CompleteQuadraticTransitionCheck, FABL.one_div_card_le_uniformProbability_of_mem, FABL.uniformProbability}
\difficulty{3}
Let $n$ be a nonnegative integer, let $w_0,\ldots,w_{n-1}$ be Boolean values, and let
$P$ assign Boolean values to product positions $X_0,\ldots,X_{n-1}$ and prefix positions
$L_0,\ldots,L_n$ and $Q_0,\ldots,Q_n$. Consider the $3n+3$ checks consisting of $L_0=0$,
$Q_0=0$, and $Q_n=1$, together with, for each $i<n$, the checks $L_i+w_i=L_{i+1}$ modulo
$2$, $X_i=L_i\land w_i$, and $Q_i+X_i=Q_{i+1}$ modulo $2$. If at least one of these
checks fails, then a uniformly chosen check fails with probability at least $1/(3n+3)$.
\end{theorem}

\begin{theorem}[Complete quadratic linear pcpp system is pcpp]
\lean{FABL.completeQuadraticLinearPCPPSystem_isPCPP}
\statementsource{Analysis-of-Boolean-Functions-by-Ryan-ODonnell}{pdf-ffbe9c7b3306-p197-b012}
\label{FABL.completeQuadraticLinearPCPPSystem_isPCPP}
\leanok
\uses{FABL.CompleteQuadraticLinearCheckHolds, FABL.IsCompleteQuadraticOne, FABL.IsFar, FABL.IsPCPP, FABL.PCPPSystem.rejectionProbability, FABL.completeQuadraticFlipNeighbor, FABL.completeQuadraticFlipNeighbor_isCompleteQuadraticOne, FABL.completeQuadraticLinearPCPPSystem, FABL.completeQuadraticLinearPCPPSystem_rejectionProbability, FABL.completeQuadraticLinearProof, FABL.completeQuadraticLinearProof_checks, FABL.isCompleteQuadraticOne_of_linearChecks, FABL.one_div_checkCount_le_completeQuadratic_failedCheckProbability, FABL.relativeHammingDist, FABL.relativeHammingDist_completeQuadraticFlipNeighbor_le_two_div, FABL.uniformProbability}
\difficulty{6}
% dataset-pipeline: stub -- prose pending build_informal_statements.py
Statement pending regeneration.
\end{theorem}

\begin{definition}[Complete quadratic linear PCPP description]
\lean{FABL.completeQuadraticLinearPCPPDescription}
\statementsource{Analysis-of-Boolean-Functions-by-Ryan-ODonnell}{pdf-ffbe9c7b3306-p197-b012}
\label{FABL.completeQuadraticLinearPCPPDescription}
\leanok
\uses{FABL.IsCompleteQuadraticOne, FABL.PCPPDescription, FABL.completeQuadraticLinearPCPPSystem, FABL.completeQuadraticLinearPCPPSystem_isPCPP}
For each positive integer $n$, there exists a probabilistically checkable proof
description for Boolean strings whose complete quadratic bit is $1$. The proof has
length $n+2(n+1)$, the verifier makes $3$ queries, and the rejection rate is $1/12$.
\end{definition}

\begin{theorem}[Proof length bound for the quadratic PCPP]
\lean{FABL.completeQuadraticLinearPCPPDescription_proofLength_le_five_mul}
\proofsource{Analysis-of-Boolean-Functions-by-Ryan-ODonnell}{pdf-ffbe9c7b3306-p197-b012}
\label{FABL.completeQuadraticLinearPCPPDescription_proofLength_le_five_mul}
\leanok
\uses{FABL.completeQuadraticLinearPCPPDescription}
\difficulty{3}
For every positive natural number $n$, the proof length of the specified PCPP
description for the complete quadratic Boolean condition is at most $5n$.
\end{theorem}

\begin{theorem}[The zero-dimensional quadratic bit]
\lean{FABL.not_isCompleteQuadraticOne_zero}
\label{FABL.not_isCompleteQuadraticOne_zero}
\leanok
\uses{FABL.IsCompleteQuadraticOne, FABL.completeQuadraticBit}
\difficulty{1}
For the unique Boolean input of length zero, the complete quadratic bit is not equal to
$1$ in the field with two elements.
\end{theorem}

\begin{theorem}[Complete quadratic linear pcpp description or zero]
\lean{FABL.completeQuadraticLinearPCPPDescription_or_zero}
\statementsource{Analysis-of-Boolean-Functions-by-Ryan-ODonnell}{pdf-ffbe9c7b3306-p197-b012}
\label{FABL.completeQuadraticLinearPCPPDescription_or_zero}
\leanok
\uses{FABL.IsCompleteQuadraticOne, FABL.PCPPDescription, FABL.completeQuadraticLinearPCPPDescription, FABL.not_isCompleteQuadraticOne_zero}
\difficulty{4}
% dataset-pipeline: stub -- prose pending build_informal_statements.py
Statement pending regeneration.
\end{theorem}

\begin{definition}[Matrix bilinear form]
\lean{FABL.matrixBilinearForm}
\statementsource{Analysis-of-Boolean-Functions-by-Ryan-ODonnell}{pdf-ffbe9c7b3306-p197-b013}
\label{FABL.matrixBilinearForm}
\leanok
\uses{FABL.F₂Cube, FABL.f₂DotProduct, FABL.𝔽₂}
Let $D$ be an $n$-by-$n$ matrix over the field with two elements, and let $x,y$ be
vectors in that field's $n$-dimensional coordinate space. The associated bilinear form
is the dot product of $x$ with $Dy$.
\end{definition}

\begin{definition}[Probability of a nonzero matrix bilinear form]
\lean{FABL.matrixBilinearNonzeroProbability}
\statementsource{Analysis-of-Boolean-Functions-by-Ryan-ODonnell}{pdf-ffbe9c7b3306-p197-b013}
\label{FABL.matrixBilinearNonzeroProbability}
\leanok
\uses{FABL.F₂Cube, FABL.matrixBilinearForm, FABL.uniformProbability, FABL.𝔽₂}
Let $D$ be an $n$-by-$n$ matrix over the field with two elements. The probability that
its bilinear form is nonzero, when $x$ and $y$ are chosen independently and uniformly
from the $n$-dimensional coordinate space over that field, is the uniform probability of
the event that the form evaluated at $(x,y)$ is nonzero.
\end{definition}

\begin{theorem}[Nonzero binary dot product probability]
\lean{FABL.uniformProbability_f₂DotProduct_ne_zero}
\leanok
\uses{FABL.F₂Cube, FABL.expect_vectorWalshCharacter, FABL.f₂DotProduct, FABL.f₂ZeroIndicator_eq, FABL.uniformProbability, FABL.vectorWalshCharacter, FABL.vectorWalshCharacter_apply}
\difficulty{5}
Let $n$ be a nonnegative integer, and let $\gamma\in\mathbb{F}_2^n$, where
$\mathbb{F}_2$ is the field with two elements. The uniform probability that the dot
product over $\mathbb{F}_2$ of $\gamma$ with a uniformly chosen $x\in\mathbb{F}_2^n$
is nonzero is $0$ if $\gamma=0$ and $1/2$ otherwise.
\end{theorem}

\begin{theorem}[Nonzero matrix image probability]
\lean{FABL.one_half_le_uniformProbability_mulVec_ne_zero}
\label{FABL.one_half_le_uniformProbability_mulVec_ne_zero}
\leanok
\uses{FABL.F₂Cube, FABL.matrixProductDifference, FABL.matrixProductVerificationPredicate_iff_difference_mulVec_eq_zero, FABL.uniformMatrixProductVerificationAcceptanceProbability, FABL.uniformMatrixProductVerificationAcceptanceProbability_le_half, FABL.uniformProbability, FABL.𝔽₂}
\difficulty{5}
Let $n\in\mathbb{N}$ and let $D$ be a nonzero $n\times n$ matrix over $\mathbb{F}_2$.
For a uniformly chosen vector $x\in\mathbb{F}_2^n$, the probability that $Dx\neq 0$ is
at least $1/2$.
\end{theorem}

\begin{theorem}[Bilinear nonzero probability identity]
\lean{FABL.matrixBilinearNonzeroProbability_eq}
\label{FABL.matrixBilinearNonzeroProbability_eq}
\leanok
\uses{FABL.F₂Cube, FABL.f₂DotProduct, FABL.matrixBilinearForm, FABL.matrixBilinearNonzeroProbability, FABL.uniformProbability, FABL.uniformProbability_f₂DotProduct_ne_zero, FABL.𝔽₂}
\difficulty{5}
Let $D$ be an $n\times n$ matrix over $\mathbb F_2$. For independent uniformly
distributed vectors $x,y\in\mathbb F_2^n$, the probability that $x\cdot(Dy)$ is nonzero
equals one half the probability that $Dy$ is nonzero for a uniformly distributed
$y\in\mathbb F_2^n$.
\end{theorem}

\begin{theorem}[Nonzero bilinear form probability]
\lean{FABL.one_fourth_le_matrixBilinearNonzeroProbability}
\label{FABL.one_fourth_le_matrixBilinearNonzeroProbability}
\leanok
\uses{FABL.matrixBilinearNonzeroProbability, FABL.matrixBilinearNonzeroProbability_eq, FABL.one_half_le_uniformProbability_mulVec_ne_zero, FABL.𝔽₂}
\difficulty{2}
Let $n\in\mathbb{N}$ and let $D$ be a nonzero $n\times n$ matrix over $\mathbb{F}_2$. If
$x$ and $y$ are chosen independently and uniformly from $\mathbb{F}_2^n$, then
$\mathbb{P}(x^{\mathsf T}Dy\neq 0)\geq 1/4$.
\end{theorem}

\begin{definition}[Rank-one self matrix]
\lean{FABL.rankOneSelfMatrix}
\statementsource{Analysis-of-Boolean-Functions-by-Ryan-ODonnell}{pdf-ffbe9c7b3306-p198-b001}
\label{FABL.rankOneSelfMatrix}
\leanok
\uses{FABL.F₂Cube, FABL.𝔽₂}
For a vector $gamma$ in the $n$-dimensional coordinate space over the field with two
elements, define the matrix whose $(i,j)$ entry is $gamma_igamma_j$.
\end{definition}

\begin{definition}[Matrix pairing with two vectors]
\lean{FABL.matrixRankOnePairing}
\statementsource{Analysis-of-Boolean-Functions-by-Ryan-ODonnell}{pdf-ffbe9c7b3306-p198-b001}
\label{FABL.matrixRankOnePairing}
\leanok
\uses{FABL.F₂Cube, FABL.𝔽₂}
Let $Gamma$ be an $n$-by-$n$ matrix over the field with two elements, and let $x,y$ be
vectors in the $n$-dimensional coordinate space over that field. Their matrix pairing is
$sum_isum_j Gamma_{ij}(x_i y_j)$.
\end{definition}

\begin{theorem}[Pairing with a rank-one matrix]
\lean{FABL.matrixRankOnePairing_rankOneSelfMatrix}
\label{FABL.matrixRankOnePairing_rankOneSelfMatrix}
\leanok
\uses{FABL.F₂Cube, FABL.f₂DotProduct, FABL.matrixRankOnePairing, FABL.rankOneSelfMatrix}
\difficulty{4}
Let $n$ be a natural number and let $\gamma,x,y\in\mathbb{F}_2^n$. The pairing of the
matrix with entries $\gamma_i\gamma_j$ against $x$ and $y$ equals $(\gamma\cdot
x)(\gamma\cdot y)$ in $\mathbb{F}_2$.
\end{theorem}

\begin{theorem}[Matrix bilinear form expansion]
\lean{FABL.matrixBilinearForm_eq_matrixRankOnePairing}
\label{FABL.matrixBilinearForm_eq_matrixRankOnePairing}
\leanok
\uses{FABL.F₂Cube, FABL.f₂DotProduct, FABL.matrixBilinearForm, FABL.matrixRankOnePairing, FABL.𝔽₂}
\difficulty{4}
Let $D$ be an $n\times n$ matrix over $\mathbb F_2$, and let $x,y\in\mathbb F_2^n$. The
bilinear form $x\cdot(Dy)$ equals $\sum_{i=1}^n\sum_{j=1}^n D_{ij}x_i y_j$.
\end{theorem}

\begin{theorem}[Pairing is linear in the matrix]
\lean{FABL.matrixRankOnePairing_sub}
\label{FABL.matrixRankOnePairing_sub}
\leanok
\uses{FABL.F₂Cube, FABL.matrixRankOnePairing, FABL.𝔽₂}
\difficulty{4}
Let $n$ be a natural number, let $A,B$ be $n\times n$ matrices over $\mathbb{F}_2$, and
let $x,y\in\mathbb{F}_2^n$. The pairing of $A-B$ against $x$ and $y$ equals the pairing
of $A$ against $x,y$ minus the pairing of $B$ against $x,y$.
\end{theorem}

\begin{theorem}[Rank one consistency iff matrix bilinear form sub eq zero]
\lean{FABL.rankOneConsistency_iff_matrixBilinearForm_sub_eq_zero}
\statementsource{Analysis-of-Boolean-Functions-by-Ryan-ODonnell}{pdf-ffbe9c7b3306-p198-b001}
\label{FABL.rankOneConsistency_iff_matrixBilinearForm_sub_eq_zero}
\leanok
\uses{FABL.F₂Cube, FABL.f₂DotProduct, FABL.matrixBilinearForm, FABL.matrixBilinearForm_eq_matrixRankOnePairing, FABL.matrixRankOnePairing, FABL.matrixRankOnePairing_rankOneSelfMatrix, FABL.matrixRankOnePairing_sub, FABL.rankOneSelfMatrix, FABL.𝔽₂}
\difficulty{3}
% dataset-pipeline: stub -- prose pending build_informal_statements.py
Statement pending regeneration.
\end{theorem}

\begin{definition}[Rank-one consistency probability]
\lean{FABL.rankOneConsistencyProbability}
\statementsource{Analysis-of-Boolean-Functions-by-Ryan-ODonnell}{pdf-ffbe9c7b3306-p198-b001}
\label{FABL.rankOneConsistencyProbability}
\leanok
\uses{FABL.F₂Cube, FABL.f₂DotProduct, FABL.matrixRankOnePairing, FABL.uniformProbability, FABL.𝔽₂}
Let $gamma$ be a vector and $Gamma$ a matrix over the field with two elements, both of
dimension $n$. Choose $x,y$ independently and uniformly from the $n$-dimensional
coordinate space. The rank-one consistency probability is the probability that
$(gammacdot x)(gammacdot y)$ equals the matrix pairing of $Gamma$ with $x,y$.
\end{definition}

\begin{theorem}[Consistency for the rank-one matrix]
\lean{FABL.rankOneConsistencyProbability_eq_one}
\label{FABL.rankOneConsistencyProbability_eq_one}
\leanok
\uses{FABL.F₂Cube, FABL.f₂DotProduct, FABL.matrixRankOnePairing, FABL.matrixRankOnePairing_rankOneSelfMatrix, FABL.rankOneConsistencyProbability, FABL.rankOneSelfMatrix, FABL.uniformProbability}
\difficulty{3}
Let $n\in\mathbb N$ and let $\gamma\in\mathbb F_2^n$. For independent uniform vectors
$x,y\in\mathbb F_2^n$, the probability that $(\gamma\cdot x)(\gamma\cdot
y)=\sum_{i,j=1}^n\gamma_i\gamma_jx_i y_j$ in $\mathbb F_2$ is $1$.
\end{theorem}

\begin{theorem}[Consistency bound for a different matrix]
\lean{FABL.rankOneConsistencyProbability_le_three_fourths}
\label{FABL.rankOneConsistencyProbability_le_three_fourths}
\leanok
\uses{FABL.F₂Cube, FABL.f₂DotProduct, FABL.matrixBilinearForm, FABL.matrixBilinearNonzeroProbability, FABL.matrixRankOnePairing, FABL.one_fourth_le_matrixBilinearNonzeroProbability, FABL.rankOneConsistencyProbability, FABL.rankOneConsistency_iff_matrixBilinearForm_sub_eq_zero, FABL.rankOneSelfMatrix, FABL.uniformProbability, FABL.𝔽₂}
\difficulty{5}
Let $n\in\mathbb N$, let $\gamma\in\mathbb F_2^n$, and let $\Gamma\in\mathbb
F_2^{n\times n}$ differ from the matrix with entries $\gamma_i\gamma_j$. For independent
uniform vectors $x,y\in\mathbb F_2^n$, the probability that $(\gamma\cdot x)(\gamma\cdot
y)=\sum_{i,j=1}^n\Gamma_{ij}x_i y_j$ in $\mathbb F_2$ is at most $3/4$.
\end{theorem}
